\documentclass[11pt]{article}
\newcommand{\DoNotLoadEpstopdf}{}
\usepackage[T1]{fontenc}
\usepackage{lmodern}
\pdfglyphtounicode{tfm:lmsy10/mapsto}{007C}
\pdfglyphtounicode{tfm:lmex10/parenleftbig}{0028}
\pdfglyphtounicode{tfm:lmex10/parenrightbig}{0029}
\pdfglyphtounicode{tfm:lmex10/vextendsingle}{23D0}
\pdfglyphtounicode{tfm:lmex10/parenleftBig}{0028}
\pdfglyphtounicode{tfm:lmex10/parenrightBig}{0029}
\pdfglyphtounicode{tfm:lmex10/parenleftbigg}{0028}
\pdfglyphtounicode{tfm:lmex10/parenrightbigg}{0029}
\pdfglyphtounicode{tfm:lmex10/braceleftbigg}{007B}
\pdfglyphtounicode{tfm:lmex10/bracerightbigg}{007D}
\pdfglyphtounicode{tfm:lmex10/parenleftBigg}{0028}
\pdfglyphtounicode{tfm:lmex10/parenrightBigg}{0029}
\pdfglyphtounicode{tfm:lmex10/summationtext}{2211}
\pdfglyphtounicode{tfm:lmex10/producttext}{220F}
\pdfglyphtounicode{tfm:lmex10/integraltext}{222B}
\pdfglyphtounicode{tfm:lmex10/summationdisplay}{2211}
\pdfglyphtounicode{tfm:lmex10/integraldisplay}{222B}
\pdfglyphtounicode{tfm:lmex10/hatwide}{0302}
\pdfglyphtounicode{tfm:lmex10/tildewide}{0303}
\pdfglyphtounicode{tfm:lmex10/bracketleftBig}{005B}
\pdfglyphtounicode{tfm:lmex10/bracketrightBig}{005D}
\pdfglyphtounicode{tfm:lmex10/radicalbig}{221A}
\pdfglyphtounicode{tfm:lmex10/radicalBig}{221A}
\pdfglyphtounicode{tfm:lmex10/parenlefttp}{239B}
\pdfglyphtounicode{tfm:lmex10/parenleftex}{239C}
\pdfglyphtounicode{tfm:lmex10/parenleftbt}{239D}
\pdfglyphtounicode{tfm:lmex10/parenrighttp}{239E}
\pdfglyphtounicode{tfm:lmex10/parenrightex}{239F}
\pdfglyphtounicode{tfm:lmex10/parenrightbt}{23A0}
\usepackage{amsmath,amssymb,amsthm,mathtools}
\usepackage{microtype}
\usepackage[margin=1.05in]{geometry}
\usepackage{enumitem}
\PassOptionsToPackage{hyphens}{url}
\usepackage[hidelinks]{hyperref}
\hypersetup{
  pdftitle={Power-law violations of Yau's nodal upper bound},
  pdfauthor={OpenAI},
  pdfsubject={A fixed smooth metric with a positive power of excess nodal volume}
}
\newcommand{\R}{\mathbb R}
\newcommand{\N}{\mathbb N}
\newcommand{\C}{\mathbb C}
\newcommand{\E}{\mathbb E}
\newcommand{\Pp}{\mathbb P}
\newcommand{\dd}{\mathrm d}
\newcommand{\Hh}{\mathcal H}
\newcommand{\Z}{\mathcal Z}
\DeclareMathOperator{\supp}{supp}
\DeclareMathOperator{\diver}{div}
\DeclareMathOperator{\tr}{tr}
\newtheorem{theorem}{Theorem}[section]
\newtheorem{proposition}[theorem]{Proposition}
\newtheorem{lemma}[theorem]{Lemma}
\newtheorem{corollary}[theorem]{Corollary}
\theoremstyle{definition}
\newtheorem{definition}[theorem]{Definition}
\theoremstyle{remark}
\newtheorem{remark}[theorem]{Remark}
\numberwithin{equation}{section}
\allowdisplaybreaks[1]
\setlist[enumerate]{itemsep=3pt,topsep=5pt}
\setlist[itemize]{itemsep=3pt,topsep=5pt}
\emergencystretch=1em
\title{Power-law violations of\\Yau's nodal upper bound}
\author{OpenAI}
\date{September 23, 2026}
\begin{document}
\maketitle
\begin{abstract}
We construct a smooth Riemannian metric on $S^4\times S^1$ and a sequence
of real Laplace eigenfunctions whose nodal four-volume grows faster than
$\lambda^{1/2+\epsilon_0}$ for one fixed $\epsilon_0>0$.
This disproves the smooth upper-bound assertion in Yau's nodal-set
conjecture and the proposed bound with an arbitrarily small positive
power loss.
\end{abstract}
\section{Introduction}\label{sec:introduction}

Let $(M^n,g)$ be a smooth closed connected Riemannian manifold. We use
$\Delta_g=\diver_g\nabla_g$ and write
\[
 -\Delta_g u=\lambda u,\qquad
 \Z(u)=\{x\in M:u(x)=0\}.
\]
Yau posed the nodal-size question in his problem collection
\cite{Yau1982}. In its smooth closed-manifold formulation, Yau's
nodal-set conjecture predicts that the $(n-1)$-dimensional
Hausdorff measure of $\Z(u)$ is bounded above and below by positive
constant multiples of $\sqrt\lambda$, with the constants depending only
on the fixed manifold and metric \cite[p.~162]{DonnellyFefferman1988}.
The upper assertion is required to
hold for every nonzero real eigenfunction, including all eigenfunctions
in a multiple eigenspace.

We disprove that upper assertion in the smooth category. In fact, the
assertion that one has an upper bound with every positive power loss,
\begin{equation}\label{eq:near-sharp-target}
 \Hh_g^{n-1}(\Z(u))
 \le C(M,g,\epsilon)\lambda^{1/2+\epsilon}
 \qquad\text{for every }\epsilon>0,
\end{equation}
fails.

\begin{theorem}\label{thm:main}
There exist a smooth Riemannian metric $g$ on $S^4\times S^1$, a number
$\epsilon_0>0$, and nonzero real smooth functions $u_j$ with positive
eigenvalues $\lambda_j$ such that
\[
 -\Delta_g u_j=\lambda_j u_j,\qquad
 \lambda_j\longrightarrow\infty,\qquad
 \frac{\Hh_g^4(\Z(u_j))}
      {\lambda_j^{1/2+\epsilon_0}}\longrightarrow\infty.
\]
The metric and $\epsilon_0$ are fixed throughout the sequence.
\end{theorem}

In particular, \eqref{eq:near-sharp-target} fails for every
$0<\epsilon\le\epsilon_0$ for the metric of Theorem~\ref{thm:main}.

The theorem concerns ordinary Laplace--Beltrami eigenfunctions on one
closed connected manifold. Taking a product with a fixed closed manifold
gives the same conclusion in each dimension greater than five. The
construction makes no assertion about the corresponding upper bound in
dimensions three or four.

In dimensions three and four, a companion paper constructs fixed smooth
metrics and real eigenfunction sequences with $\lambda\to\infty$ for which
the nodal measure divided by $\sqrt\lambda$ tends to infinity, without
asserting a fixed positive power-law excess
\cite[Theorems~1.1 and~1.2]{OpenAINodalQualitative2026}.

\paragraph{Background.}
Donnelly and Fefferman proved the sharp two-sided nodal estimate for
real-analytic metrics \cite[Theorem~1.2]{DonnellyFefferman1988}.
For smooth metrics, Logunov proved the sharp lower bound
\cite[Theorem~1.1]{LogunovLower2018} and, in dimensions at least three,
a polynomial upper bound with a dimension-dependent exponent
\cite[Theorem~6.2]{LogunovUpper2018}.
Regularity assumptions between analyticity and smoothness also yield
stronger upper estimates: Hezari proved a $\lambda^{s/2}$ bound in
Gevrey class $s\ge1$, and a $\sqrt\lambda\log(\lambda+e)$ bound for a
specified quasianalytic class
\cite{Hezari2023}.
Those hypotheses do not cover arbitrary smooth metrics.
Theorem~\ref{thm:main} disproves the smooth upper-bound part of Yau's
conjecture and even the relaxed estimate \eqref{eq:near-sharp-target};
the lower-bound results remain compatible with it.

There is a complementary tradition of prescribing nodal geometry.
In dimensions at least three, Enciso and Peralta-Salas realize a smooth
closed connected orientable separating hypersurface as the nodal set of
a first nonconstant eigenfunction by choosing a metric
\cite{EncisoPeraltaSalas2015}. At high energy, Enciso, Peralta-Salas,
and Torres de Lizaur realize prescribed stable nodal components, up to
ambient diffeomorphism, on round spheres and flat tori, with a
localization restriction in the torus case
\cite{EncisoPeraltaSalasTorres2018}. Their construction already gives
arbitrarily high frequencies on fixed analytic metrics. The distinction
here is power-law excess in total nodal measure, which is not
furnished by those prescription theorems.

\paragraph{The quantitative mechanism.}
We first work with a symmetric elliptic operator on the four-sphere
whose positive principal tensor and positive density can be varied
independently. A local corrugation increases the integral of the
gradient of a logarithmic envelope. More specifically, the part of a
parent cube prepared for another corrugation has more new gradient
integral than the entire old parent cube. Repetition at a fixed scale
ratio therefore gives a multiplicative gain. The derivatives of order
$r$ grow at most exponentially in $(r+1)$ times the number of
generations. Let $k$ denote the construction frequency, with target
eigenvalue $k^2$. Taking a fixed small multiple of $\log k$ generations
produces a smooth real envelope $h=h_k$ on the base and a fixed
coordinate region $V$ such that
\[
 \int_V\sqrt{1+|\nabla h|^2}\,\dd x\ge c k^\gamma
\]
for one fixed $\gamma>0$, while keeping every fixed derivative of $h$
within a controlled power of $k$.

Finite Taylor solutions of the complex eikonal and transport equations
then give localized packets beneath $e^{kh}$. They achieve any prescribed
finite residual accuracy while preserving fixed derivative exponents.
A real Gaussian combination has a quantitatively nonvanishing first
jet everywhere and opposite signs on many pairs of small disks.
The first property permits an explicit simultaneous correction of the
principal tensor and density to make the function an exact
eigenfunction. The second gives stable lower bounds for nodal measure
without a regularity assumption on the zero set.

The distinction between constants and exponents is essential here.
At a given stage, constants and frequency thresholds may depend
arbitrarily badly on the required finite coefficient accuracy.
The exponent $\gamma$ does not; nor do the chart or coefficient
neighborhood. The positive constant in the nodal lower bound and the frequency threshold
may depend on the fixed input pair, collar, finite norm, and tolerance.
Consequently the construction can
meet every finite collection of earlier spectral-persistence tolerances
and still retain a fixed power gain. A summable iteration gives one
smooth limiting pair. Adding a warped circle realizes its weighted
operator as an ordinary Laplace--Beltrami operator.

These steps also give reusable local facts: a corrugation with a
repeatable integral gain, a phase-Hessian construction valid at zero
real slope, a polynomial first-jet noncancellation argument, and an
exact tensor--density correction. Each is proved below. In particular,
the argument does not appeal to a qualitative existence theorem with an
uncontrolled frequency threshold.

\paragraph{Organization.}
Section~\ref{sec:framework} states the quantitative one-step construction
and fixes the geometric notation.
Section~\ref{sec:envelopes} proves the envelope estimates.
Section~\ref{sec:waves} constructs and estimates the local packets.
Section~\ref{sec:global} produces exact eigenfunctions, proves the
one-step construction, and passes to the fixed metric in
Theorem~\ref{thm:main}.

\section{Framework and the quantitative step}\label{sec:framework}

\subsection{The weighted base operator}

Let $B=S^4\subset\R^5$, with round metric $g_0$ and round Riemannian
density $\mu$. Write a point of $B$ as $(x_1,x_2,y)$, where $y\in\R^3$,
and set $t=|y|^2$. A \emph{positive pair} consists of a smooth positive
definite symmetric contravariant tensor $A$ and a smooth function
$\rho>0$. Define a metric $G$ and a differential operator $L$ by
\begin{equation}\label{eq:weighted-operator}
 G^{-1}=\rho^{-1}A,\qquad
 Lf=\rho^{-1}\diver_\mu(A\,\dd f).
\end{equation}
Thus $-L$ is symmetric and nonnegative in $L^2(B,\rho\,\dd\mu)$, with
quadratic form
\[
 \mathfrak a(f,f)=\int_B A(\dd f,\dd f)\,\dd\mu .
\]
The tensor and density will be changed independently. Their conversion
to an ordinary metric occurs only at the end of the proof.

Fix a finite smooth atlas with compact coordinate subpatches covering
$B$; all $C^r$ norms of functions and tensors may be taken in that atlas.
For a pair, the norm means the sum of the two component norms.
Equivalent fixed choices change constants but do not change the
assertions. All pairs used below lie in one sufficiently small $C^1$
neighborhood of $(g_0^{-1},1)$, within the space of smooth positive
pairs. Accordingly $G$ and $G^{-1}$ have common $C^1$ bounds and common
ellipticity constants. Higher derivative bounds need only be finite for
each fixed input pair.

Gradients, Hessians, lengths, and orthogonality in envelope and phase
estimates refer to $G$ unless another metric is displayed. Integration
in the elliptic operator refers to $\mu$; integration in a coordinate
corrugation chart refers to Lebesgue measure $\dd x$.
For a real smooth function $h$, set
\begin{equation}\label{eq:hessian-test}
 \begin{split}
 p&=\nabla_G h,\qquad Q=(1+|p|_G^2)^{1/2},\\
 J_h(e)&=(\nabla_G^2h)(p/Q,p/Q)+(\nabla_G^2h)(e,e),
 \qquad |e|_G=1,\quad \dd h(e)=0 .
 \end{split}
\end{equation}
At $p=0$ every unit vector is an admissible test.
The positivity conditions imposed below on $J_h$ supply a gap between
the envelope Hessian and the real part of a complex phase Hessian.
Lemma~\ref{lem:phase-hessian} and Proposition~\ref{prop:packets} use
this gap in the construction and localization of the packets.
Complex inner products
and tensor evaluations in the phase equations are extended
bilinearly; estimates of the size of a complex vector or matrix use
the ordinary Hermitian norm.

We fix once and for all
\[
 \delta=10^{-6}.
\]
The numerical exponents below leave generous margins; we do not optimize
the final positive gain.
Positive constants written $c,C,C_r$ may change between occurrences.
Any additional dependencies are specified where their uniformity
matters. In particular, a finite approximation order is fixed before a
frequency tends to infinity.

\subsection{A round region and sign certificates}

For an integer $N>1$, let
\begin{equation}\label{eq:round-beam}
 k=\sqrt{N(N+3)},\qquad
 u_*=\operatorname{Re}(x_1+i x_2)^N .
\end{equation}
The Euclidean polynomial in \eqref{eq:round-beam} is homogeneous of
degree $N$ and harmonic in $\R^5$. The polar-coordinate formula
\[
 \Delta_{\R^5}
 =\partial_r^2+\frac4r\partial_r+\frac1{r^2}\Delta_{g_0}
\]
therefore gives $-\Delta_{g_0}u_*=k^2u_*$ on $B$.
Where $t<1$, the angular coordinate $\vartheta$ in the
$(x_1,x_2)$-plane also gives
\[
 u_*=(1-t)^{N/2}\cos(N\vartheta).
\]
The round beam will be retained in part of a region where the input
pair is round. This region serves two purposes: it supplies an exact
solution to which the packets can be glued, and it leaves a patch where
a later tensor perturbation can make the eigenvalue simple while
preserving that solution. The construction uses both properties before
passing to the limiting operator.

The envelope construction fixes a coordinate chart
$\Omega\subset\{t>1/2\}$ and an open region $V$ whose closure is compact
in $\Omega$. The initial corrugation cube is compactly contained in
$V$. Euclidean balls centered in $V$ with radii tending uniformly to
zero are consequently contained in $\Omega$.

\begin{definition}\label{def:certificate}
A \emph{sign certificate} for a continuous real function $u$ in
$\Omega$ is a finite family of pairs $(D_i^-,D_i^+)$ of congruent
closed parallel three-dimensional Euclidean disks such that:
\begin{enumerate}[label=\textup{(\roman*)}]
 \item each pair is the two ends of a straight closed cylinder
 compactly contained in a Euclidean coordinate ball;
 \item these coordinate balls are pairwise disjoint;
 \item $u$ has strictly opposite signs on $D_i^-$ and $D_i^+$.
\end{enumerate}
The \emph{size} of the certificate is
$\sum_i\Hh_{\mathrm{Euc}}^3(D_i^+)$.
\end{definition}

The signs in Definition~\ref{def:certificate} can be reversed on each
pair. Since the disks are compact and finite in number, a certificate
persists under sufficiently small uniform perturbations of $u$.
It is also preserved by multiplying $u$ by any nonzero real constant.
Every segment joining corresponding points of $D_i^-$ and $D_i^+$
contains a zero. This elementary projection property will provide the
Hausdorff-measure bound for the final metric.

\subsection{The one-step construction}

\begin{proposition}\label{prop:step}
There exist a $C^1$ neighborhood $\mathcal U$ of the round pair
$(g_0^{-1},1)$ among smooth positive pairs, a chart $\Omega$ as above,
and a number $\gamma>0$ with the following property.
Fix a smooth pair $(A,\rho)\in\mathcal U$, numbers
\[
 \frac14\le l_-<l_+\le\frac12,
\]
and suppose that $(A,\rho)=(g_0^{-1},1)$ on $\{t<l_+\}$.
For every integer $r\ge0$ and every $\eta>0$, there exist $c>0$ and
$N_0$ such that for every integer $N\ge N_0$, with $k$ as in
\eqref{eq:round-beam}, there are a smooth positive pair
$(\widehat A,\widehat\rho)\in\mathcal U$ and a nonzero real smooth
function $u$ satisfying:
\begin{enumerate}[label=\textup{(\roman*)}]
 \item $\|(\widehat A-A,\widehat\rho-\rho)\|_{C^r}<\eta$, and
 $(\widehat A,\widehat\rho)=(g_0^{-1},1)$ on $\{t<l_-\}$;
 \item for $\widehat L$ defined by \eqref{eq:weighted-operator},
 $-\widehat Lu=k^2u$, and $k^2$ is a simple eigenvalue of
 $-\widehat L$ on the base $B$;
 \item $u$ has a sign certificate in $\Omega$ of size at least
 $c k^{1+\gamma}$.
\end{enumerate}
The neighborhood $\mathcal U$, chart $\Omega$, and exponent $\gamma$
are common to all stages. The constant $c$ and threshold $N_0$ may
depend on the fixed input pair, $l_-,l_+,r,\eta$.
\end{proposition}

The proof of Proposition~\ref{prop:step} occupies the next three
sections. Membership of the output in $\mathcal U$ is ensured by
imposing an additional sufficiently small tolerance around the fixed
input pair. The conclusion permits arbitrarily small finite-norm
tolerances; it makes no claim of a uniform frequency threshold as
those tolerances or the collar width vary. The common exponent
$\gamma$ is the feature that allows a smooth limiting metric to retain
a fixed power of excess nodal growth.

\section{Envelopes with a fixed power gain}\label{sec:envelopes}

We construct envelopes whose gradient integral grows geometrically under
iteration.  The part of each parent cube available for the next iteration
will have more gradient integral than the entire parent had before the
change.  All scales and gain factors are fixed before any finite
approximation order is prescribed.

Throughout this Section, tensor norms in Hessian inequalities are operator
norms for $G$.  Coordinate derivatives and integrals refer to fixed smooth
coordinates.  The notation $p,Q,J_h$ is that of
\eqref{eq:hessian-test}.

\subsection{A seed and the preparation condition}

\begin{lemma}\label{lem:seed}
There are a number $a_*>0$, a smooth real function $h_0$ on $B$, a
$C^1$ neighborhood $\mathcal G$ of $g_0$ among smooth metrics, and a
coordinate chart $\Omega\subset\{t>1/2\}$ with open sets
$U_0\Subset V\Subset\Omega$ such that the following assertions hold
uniformly for $G\in\mathcal G$:
\begin{enumerate}[label=\textup{(\roman*)}]
 \item $h_0=-a_*t$ on $\{t\le1/2\}$;
 \item $\max_e J_{h_0}(e)\ge c_0>0$ on $\{t\ge1/4\}$;
 \item on $U_0$, the Hessian of $h_0$ is positive definite with a
 uniform positive lower bound, and
 $0<s_{\min}\le|\nabla_Gh_0|_G\le s_{\max}$.
\end{enumerate}
The maximum in \textup{(ii)} is over all admissible tests, including all
unit vectors when $\dd h_0=0$.
\end{lemma}

\begin{proof}
Begin with $\widetilde h=-a_*t$.  Direct differentiation on the round
sphere gives
\begin{equation}\label{eq:seed-identities}
 |\nabla_{g_0}\widetilde h|_{g_0}^2
 =4a_*^2t(1-t),\qquad
 \nabla_{g_0}^2\widetilde h
 =2a_*(tg_0-\dd y\cdot\dd y).
\end{equation}
At every point there is a nonzero tangent vector in $\ker\dd y$,
because $\dd y:T_xB\to\R^3$ has a four-dimensional domain.  A round
unit vector $e$ in this kernel satisfies $\dd\widetilde h(e)=0$ and
$(\nabla_{g_0}^2\widetilde h)(e,e)=2a_*t$.  Since the Hessian in
\eqref{eq:seed-identities} has norm at most $2a_*$ and the squared
gradient has norm at most $a_*^2$, on $t\ge1/4$ we obtain
\[
 J_{\widetilde h}(e)\ge a_*/2-2a_*^3\ge a_*/4
\]
after choosing $a_*^2\le1/8$.

At a point $x_*$ with $t=1$, the Hessian has eigenvalues
$2a_*,2a_*,0,0$, and the gradient is zero.  On a small neighborhood
of $x_*$ contained in $\{t>1/2\}$, two Hessian eigenvalues therefore
have a uniform positive lower bound and the gradient is arbitrarily
small.  In round normal coordinates $\xi$ at $x_*$, add
$\varepsilon|\xi|^2/2$ times a fixed smooth cutoff equal to one near
$\xi=0$.  Choose $\varepsilon>0$ sufficiently small that this is a
$C^2$-small perturbation on its support.  The resulting function $h_0$
has positive definite Hessian at $x_*$.

Here is why the cutoff region retains the required test.  Shrinking its
support and then the perturbation if necessary, the perturbed Hessian
has two eigenvalues at least $c a_*$, its norm is at most $C a_*$,
and $|\nabla_{g_0}h_0|$ is sufficiently small that
$C a_*|\nabla_{g_0}h_0|^2<c a_*/2$.  The positive two-plane
intersects $\ker\dd h_0$ in dimension at least one.  A unit vector
in that intersection gives $J_{h_0}\ge c a_*/2$.  This argument also
applies at a critical point.  Outside the perturbation support the
previous $\ker\dd y$ test applies.

These conclusions persist in a sufficiently small common $C^1$
neighborhood of $g_0$.  Indeed, the difference of the two Hessians of
a fixed function is the connection difference contracted with its
differential, and all metric and gradient comparisons use only
$C^1$ bounds.  Outside the perturbation support, the condition
$\dd y(e)=0$ still implies $\dd h_0(e)=0$ independently of the
metric.  On the support, the positive two-plane and small-gradient
estimates just established are uniform under a small perturbation.

Finally, positive definiteness at $x_*$ holds on a neighborhood of that
point.  There is a nearby point with nonzero differential: otherwise
$h_0$ would be constant on that neighborhood.  Choose a relatively
compact neighborhood $U_0$ of such a point where the Hessian remains
positive definite and the gradient stays separated from zero.  Choose
$U_0\Subset V\Subset\Omega\subset\{t>1/2\}$ in a fixed
coordinate chart, and shrink $\mathcal G$ once more.  All the stated
constants can then be chosen uniformly.  Notice that $U_0$ is chosen
away from the critical point $x_*$.
\end{proof}

We use the coordinates of $\Omega$ for all corrugations.  They may be
taken on a neighborhood of its closure, so metric and inverse metric
have common ellipticity and $C^1$ bounds there.  Fix positive constants
$C_{\mathrm p}$ and $c_{\mathrm p}$.  A smooth function $h$ is
\emph{prepared at scale $S$} on a cube $P$ of side $\nu S$,
$0<S\le1$, if
\begin{equation}\label{eq:preparation}
 \begin{gathered}
 |\nabla h|>0,\qquad
 |\nabla^2h|\le C_{\mathrm p}|\nabla h|/S,\\
 \text{at each point there is a two-plane }
 E\subset\ker\dd h\text{ such that }
 J_h(e)\ge c_{\mathrm p}|\nabla h|/S
 \quad(e\in E,\ |e|=1).
 \end{gathered}
\end{equation}
The plane $E$ need not vary continuously.  Cubes are parallel to the
fixed coordinate axes, and the inequalities are required on their
closures.

The two Hessian conditions serve different parts of the next corrugation.
The Hessian norm bound controls the inherited Hessian and the variation
of the parent gradient.  When the gradient changes by a small relative
amount, we can rotate the favorable two-plane slightly into the
hyperplane perpendicular to the new gradient.  One further linear
tangency constraint then leaves a direction with a positive test.
On the region selected for repetition, the modification supplies two
favorable directions again.  This is why the preparation condition
asks for a two-plane even though the envelope ultimately needs only one
positive test at each point.

By Lemma~\ref{lem:seed}, for all sufficiently large
$C_{\mathrm p}$ and sufficiently small $c_{\mathrm p}>0$, every
sufficiently small cube in $U_0$ is prepared for $h_0$ at scale $1$,
uniformly over $G\in\mathcal G$.

\subsection{A repeatable local corrugation}

\begin{lemma}\label{lem:corrugation}
Under fixed ellipticity and $C^1$ metric bounds in the coordinate
patch, there are constants $C_{\mathrm p}^*>0$,
$c_{\mathrm p}^*>0$ and $\zeta>0$ with the following property.
For any $C_{\mathrm p}\ge C_{\mathrm p}^*$ and
$0<c_{\mathrm p}\le c_{\mathrm p}^*$, one may choose a sufficiently
small reciprocal integer $\nu>0$, once and for all, such that every
prepared cube $P$ at scale $S\le1$ admits a smooth modification
$v=h+w$, with $w$ supported in the interior of $P$, satisfying:
\begin{enumerate}[label=\textup{(\roman*)}]
 \item for fixed constants $0<\alpha\le1\le\beta$ and $c>0$,
 \[
  \alpha|\nabla h|\le|\nabla v|\le\beta|\nabla h|,
  \qquad
  \max_eJ_v(e)\ge c|\nabla h|/S
  \quad\hbox{on }P;
 \]
 \item writing $\ell=\nu^2S$, some cubes in the partition of $P$
 into cubes of side $\nu\ell$ are prepared for $v$ at scale $\ell$
 with the same constants in \eqref{eq:preparation}, and their union
 $\mathcal R_P$ satisfies
 \begin{equation}\label{eq:local-mass-gain}
  \int_{\mathcal R_P}|\nabla v|\,\dd x
  \ge(1+\zeta)\int_P|\nabla h|\,\dd x;
 \end{equation}
 \item if $s=|\nabla h(x_P)|$ at the center $x_P$ of $P$, then for
 every integer $j\ge0$,
 \begin{equation}\label{eq:corrugation-derivatives}
  \|D^jw\|_{L^\infty(P)}\le C_j s\ell^{1-j}.
 \end{equation}
\end{enumerate}
All constants are uniform over prepared functions and cubes and over
the stated class of metrics.  They may depend on the fixed pattern,
preparation constants, and final choice of $\nu$.  No derivative of
the input $h$ beyond its prepared second-derivative bound is used.
The pattern and $\zeta$ are chosen before the preparation constants
and the final $\nu$; increasing $C_{\mathrm p}$ or decreasing
$c_{\mathrm p}$ only requires taking $\nu$ smaller.
\end{lemma}

\begin{proof}
The pattern will depend on three variables transverse to the parent
gradient frozen at its center.  Radial wells in these variables have
two positive tangential Hessian directions; a small amplitude controls
their potentially negative radial contribution to $J_v$.  We first
choose the wells and smaller concentric balls that will supply the
repeatable gain.  We then place this pattern in the parent and verify
the Hessian test, including where the parent cutoff varies.

\paragraph{The periodic wells.}
Choose a smooth function $F$ on $[0,1]$, positive on $(0,1)$, equal
to $v$ near $v=0$, and flat at $v=1$.  Inside a unit cube in $\R^3$
take finitely many balls with pairwise disjoint closures and radii
$R_i>0$.  On a ball, at radius $r$ from its center, set
\[
 \psi(r)=-\int_r^{R_i}F(q/R_i)\,\dd q,\qquad
 f(r)=F(r/R_i),\qquad u_0(r)=f(r)/r.
\]
Extend $\psi$ by zero outside the balls and periodically to $\R^3$.
It is smooth: near a center it is a constant plus $r^2/(2R_i)$,
and it is flat at the outer boundary.  The extended value of $u_0$
at the center is $1/R_i$.  Write $\sigma$ for the radial unit vector.
Then
\begin{equation}\label{eq:well-identities}
 \nabla\psi=f\sigma,\qquad
 D^2\psi=f'\sigma\otimes\sigma
              +u_0(I-\sigma\otimes\sigma),\qquad
 f^2|f'|\le C u_0,\quad f^2\le C u_0.
\end{equation}
The first constant in these inequalities is independent of the radii,
because, for $v=r/R_i$,
\[
 \frac{f^2|f'|}{u_0}=vF(v)|F'(v)|.
\]
The second follows from $f^2/u_0=R_i vF(v)$ and the bounded radii.
Set $f,f',u_0$ equal to zero outside the balls.  These coefficients
are continuous, and
\begin{equation}\label{eq:well-vanishing}
 u_0=0\quad\Longrightarrow\quad f=f'=\psi=0.
\end{equation}
At a center $f'=u_0$, so any choice of $\sigma$ gives the same
Hessian in \eqref{eq:well-identities}; outside the balls choose
$\sigma$ arbitrarily as well.

Choose a small amplitude $a>0$.  Its smallness will use only $F$ and
the radius-independent first bound in \eqref{eq:well-identities}.
In particular it can be chosen before the packing.  The mean gradient
factor on any one well is
\[
 \mathcal W(a)=3\int_0^1v^2\sqrt{1+a^2F(v)^2}\,\dd v>1.
\]
Put $\eta_0=(\mathcal W(a)-1)/8$.  Choose the finite packing with
volume fraction $\varpi$ such that
\[
 \varpi\mathcal W(a)>1+6\eta_0.
\]
Such a packing exists.  To see this without requiring an infinite
packing in the construction, fill a fixed positive fraction of any
remaining open set by finitely many closed grid cubes lying inside
it, and put a ball with strictly smaller radius than the inradius
in each cube.  This captures a fixed positive fraction of the
remaining volume.  Finitely many repetitions make the uncovered
volume as small as desired, with all ball closures still disjoint.

Choose $\eta>0$ so small that, on the core balls
$r<(1-\eta)R_i$, the weighted integral per period satisfies
\[
 C_{\mathrm{core}}
 :=\sum_i\int_{r<(1-\eta)R_i}\sqrt{1+a^2f^2}\,\dd z'
 >\varpi\mathcal W(a)-\eta_0>1+5\eta_0.
\]
Choose an inner cube fraction $q\in(0,1)$ such that
\begin{equation}\label{eq:repeat-slack}
 qC_{\mathrm{core}}>C_{\mathrm{core}}-\eta_0>1+4\eta_0.
\end{equation}
Here $q$ is a four-dimensional volume fraction.  Fix a smooth cutoff
on the unit parent cube which equals one on the concentric inner
cube of fraction $q$, is supported in the parent interior, and is
flat wherever it is zero.  Products of one-dimensional plateau
cutoffs give such a function.  The wells, cores, and cutoff are now
fixed.  All their derivative bounds are finite, although no bound
uniform in derivative order is asserted.

\paragraph{Normalization on a prepared parent.}
Take affine coordinates $z=(z_1,z')$ centered at $x_P$, making the
metric Euclidean at the center and
$\nabla h(x_P)=s\partial_{z_1}$.  The linear maps and their inverses
have uniform bounds; the image of $P$ is a parallelepiped of uniformly
bounded distortion.  The coordinate identity relating $D(\dd h)$
to the covariant Hessian and the bounded connection gives
\[
 |D(\dd h)|\le C|\dd h|/S
\]
on $P$, since $S\le1$.  Integration along segments of length
$O(\nu S)$ yields relative upper and lower bounds for the gradient
and then
\begin{equation}\label{eq:parent-normalization}
 \dd h/s=\dd z_1+O(\nu),\qquad
 (G_{ij})_z=I+O(\nu),\qquad
 |\nabla h|/s=1+O(\nu).
\end{equation}
The constants here may depend on the preparation constants.  Only
the bounds already specified in the lemma enter them.

Rescale the fixed parent cutoff to $P$ and call it $\chi$.  Set
\begin{equation}\label{eq:corrugation-formula}
 v=h+s\ell a\chi\psi(z'/\ell),\qquad \ell=\nu^2S.
\end{equation}
In the following estimates the radial coefficients are evaluated at
$z'/\ell$.  A cutoff derivative at its parent scale is denoted
$D_{\mathrm{sc}}\chi$; thus an ordinary derivative of $\chi$ is
$(\nu S)^{-1}D_{\mathrm{sc}}\chi$.  The bounded affine changes only
alter the fixed bounds for these scaled derivatives.  Since
$\ell/(\nu S)=\nu$, Equation~\eqref{eq:corrugation-formula} gives
\begin{equation}\label{eq:new-gradient}
 \dd v/s=\dd z_1+a\chi f\sigma\cdot\dd z'+O(\nu).
\end{equation}
This proves the gradient comparison in \textup{(i)} for fixed
$\alpha,\beta$, when $\nu$ is sufficiently small.

\paragraph{The new Hessian contribution.}
Write
\[
 D=\frac{\nabla v}{\sqrt{1+|\nabla v|^2}},\qquad
 \mathcal F=\{e:\dd v(e)=0,\ \sigma\cdot e'=0\},
\]
where $e'$ denotes the static transverse components in the affine
coordinates.  The two defining linear conditions are independent
by \eqref{eq:new-gradient}, so $\mathcal F$ is two-dimensional.
For a metric unit vector $e\in\mathcal F$,
\begin{equation}\label{eq:test-components}
 |e'|^2\ge\tfrac12,\qquad
 |\sigma\cdot D'|\le C_0a\chi f+C\nu.
\end{equation}
Indeed $e_1=O(\nu)$ follows from the two linear conditions.  The
second estimate follows from \eqref{eq:new-gradient} and
$G^{ij}=\delta_{ij}+O(\nu)$, using a denominator at least a fixed
multiple of $s$.  We may take $C_0$ absolute by first taking $\nu$
small for the fixed data.  These arguments are uniform even when
$s$ is arbitrarily small or large.

Let $\Pi_\sigma$ be Euclidean projection onto $\sigma^\perp$ in
the transverse space.  The leading contribution of $v-h$ to its
Hessian test on $D,e$, divided by $s/\ell$, is
\[
 a\chi\left[
 f'(\sigma\cdot D')^2+
 u_0\bigl(|\Pi_\sigma D'|^2+|e'|^2\bigr)\right].
\]
By \eqref{eq:test-components}, this is at least
\[
 a\chi\left[
 \tfrac12u_0-|f'|\bigl(2C_0^2a^2\chi^2 f^2+C\nu^2\bigr)
 \right].
\]
At the initial choice of amplitude, impose
$2C_0^2a^2\sup_v vF(v)|F'(v)|\le1/4$.
The main term is then at least
$a\chi u_0/4-Ca\nu^2\chi|f'|$.

The mixed cutoff terms on $e,e$ vanish because
$\sigma\cdot e'=0$.  On $D,D$ their size is bounded by
\[
 C\nu f|D_{\mathrm{sc}}\chi|(\chi f+\nu).
\]
The pure cutoff terms have size at most
$C\nu^2|\psi||D_{\mathrm{sc}}^2\chi|$.
Finally, replacing the coordinate Hessian of $v-h$ by its covariant
Hessian gives an error of size
\[
 C\ell(\chi f+\nu)
\]
in the same units, because the connection is uniformly bounded.

We verify uniformity of these errors also where the favorable
coefficient vanishes.  Set $X=\chi u_0$.  On a fixed compact
parameter space, if continuous functions $X,b\ge0$ satisfy
$b=0$ on $\{X=0\}$, then
\begin{equation}\label{eq:vanishing-coefficient}
 \sup\frac{\nu^2b}{X+\nu^2}\longrightarrow0
 \qquad(\nu\longrightarrow0).
\end{equation}
For $X\ge\tau$ the quotient is at most
$\nu^2\sup b/\tau$; for $X<\tau$ it is at most
$\sup_{X<\tau}b$, which tends to zero as $\tau\downarrow0$.
The compact parameter space here consists of cutoff position,
period position, and bounded affine frames.  Taking its closure is
harmless.  Transformed cutoff derivatives are continuous on this
space and have uniform bounds.

Apply \eqref{eq:vanishing-coefficient} to
$b=\chi|f'|$, $b=f|D_{\mathrm{sc}}\chi|$, and
$b=|\psi||D_{\mathrm{sc}}^2\chi|$.  Their vanishing follows from
\eqref{eq:well-vanishing} and the flatness of $\chi$ on its zero
set.  The other mixed term satisfies
$\nu\chi f^2|D_{\mathrm{sc}}\chi|\le C\nu X$.
Moreover, $\chi f\le C\sqrt X$ and $\ell\le\nu^2$ give
\[
 \frac{\ell(\chi f+\nu)}{X+\nu^2}
 \le C\frac{\nu^2\sqrt X}{X+\nu^2}+\nu
 \le C'\nu.
\]
We have proved, uniformly in all prepared parents and allowed tests,
\begin{equation}\label{eq:corrugation-hessian-error}
 \frac{\ell}{s}\left[
  \nabla^2(v-h)(D,D)+\nabla^2(v-h)(e,e)\right]
 \ge c_a X-\varepsilon(\nu)(X+\nu^2),
 \qquad \varepsilon(\nu)\longrightarrow0,
\end{equation}
for a fixed $c_a>0$.  No positive lower bound on a nonzero value of
$\chi$ or $u_0$ has been used.

\paragraph{Retaining one direction everywhere.}
The old Hessian has norm $O(\nu^2)$ after division by $s/\ell$.
Thus the preceding estimate compares a favorable new contribution of
size $X$ with inherited terms of size $\nu^2$.  When $X$ dominates
$\nu^2$, the new favorable directions suffice.  When $X=O(\nu^2)$,
the gradients are close enough to carry an old favorable direction to
the new tangent space.  To obtain one uniform lower bound, consider
any sequence $\nu\to0$ of admissible configurations.
If $X/\nu^2\to\infty$ along a subsequence, the positive term in
\eqref{eq:corrugation-hessian-error} dominates the old Hessian in
every unit direction of $\mathcal F$.

If instead $X/\nu^2$ stays bounded, then
$\chi f\le C\sqrt X=O(\nu)$, so
$|\nabla v-\nabla h|\le\varepsilon_1|\nabla h|$ with
$\varepsilon_1=O(\nu)$.  This also makes their normalized gradients
close uniformly for unbounded slopes.  In fact, for
$T(p)=p/\sqrt{1+|p|^2}$ and $|q-p|\le\varepsilon_1|p|$,
$\varepsilon_1\le1/2$, the segment from $p$ to $q$ stays outside
the ball of radius $|p|/2$.  Therefore
\begin{equation}\label{eq:normalized-gradient-stability}
 |T(q)-T(p)|
 \le\frac{\varepsilon_1|p|}{\sqrt{1+|p|^2/4}}
 \le2\varepsilon_1.
\end{equation}
There is a metric orthogonal rotation $R=I+O(\varepsilon_1)$ taking
$\ker\dd h$ to $\ker\dd v$.  Rotate the old good two-plane $E$
in \eqref{eq:preparation} by $R$.  Inside the three-dimensional
space $\ker\dd v$, the planes $RE$ and $\mathcal F$ intersect in
dimension at least
\[
 2+2-3=1.
\]
For a unit vector $e$ in their intersection, compare the old test on
$p/Q,R^{-1}e$ with the old Hessian evaluated on $D,e$.
Equation~\eqref{eq:normalized-gradient-stability} shows that the
change costs $o(\nu^2)$ in the normalized units.  Thus the old
contribution is at least $\tfrac12c_{\mathrm p}\nu^2$ for small
$\nu$.  The new contribution in
\eqref{eq:corrugation-hessian-error} is bounded below by
$-o(\nu^2)$ in this bounded-ratio regime.

These alternatives prove the uniform margin
$c_{\mathrm p}\nu^2/8$ in the normalized units for some test when
$\nu$ is sufficiently small.  Indeed, a contrary
sequence has a subsequence on which $X/\nu^2$ is bounded or tends
to infinity, and both alternatives just excluded failure of that
margin.  Multiplication by $s/\ell$ and
\eqref{eq:parent-normalization} prove the Hessian assertion in
\textup{(i)}.  This is where base dimension four is used: in
dimension $d$ the corresponding intersection lower bound is $d-3$.

\paragraph{The region prepared for repetition.}
Inside the inner parent cube where $\chi=1$, keep the points for
which $z'/\ell$ lies in a core ball.  On this region $u_0$ has a
fixed positive lower bound.  Equations
\eqref{eq:corrugation-hessian-error} and
\eqref{eq:parent-normalization} imply, for every unit
$e\in\mathcal F$,
\[
 J_v(e)\ge c' s/\ell,
 \qquad
 |\nabla^2v|\le(C'+O(C_{\mathrm p}\nu^2))s/\ell.
\]
The leading constants depend only on the fixed pattern and the
uniform metric bounds.  Gradient comparison now permits fixed
thresholds $C_{\mathrm p}^*$ and $c_{\mathrm p}^*$ such that, for
$C_{\mathrm p}\ge C_{\mathrm p}^*$ and
$0<c_{\mathrm p}\le c_{\mathrm p}^*$, this region satisfies
\eqref{eq:preparation} at scale $\ell$ once $\nu$ is small.
In particular the old term $O(C_{\mathrm p}\nu^2)$ is absorbed
after $C_{\mathrm p}$ has been fixed.

Select all closed child cubes contained in this open repeat region.
Since $\nu$ is a reciprocal integer, the parent has an exact
partition into cubes of side $\nu\ell$.  These children satisfy the
preparation condition.  It remains to show that their total gradient
integral exceeds that of the entire old parent.  To evaluate their mass,
first tile the affine $z$ coordinates by period cells of side
$\ell$, using a dummy period in $z_1$.  The child grid uses the original
coordinates, whereas these period cells use the affine coordinates;
the two grids need not align.  The cells meeting the inner
parent boundary have relative volume
$O(\ell/(\nu S))=O(\nu)$, uniformly over the bounded affine
distortions.  The points lost by retaining only whole child cubes
lie in a layer of width $O(\nu\ell)$ along the core and inner
parent boundaries.  In each rescaled period the core boundaries
have finite area, so this gives another relative $O(\nu)$ loss.
The constants may involve the smallest radius in the finite fixed
packing.

On the repeat region,
\[
 |\nabla v|/s=\sqrt{1+a^2f^2}+O(\nu).
\]
The constant affine Jacobian cancels on dividing by the parent
volume.  Consequently
\[
 \frac1{s|P|}\int_{\mathcal R_P}|\nabla v|\,\dd x
 \ge qC_{\mathrm{core}}-C\nu,
 \qquad
 \frac1{s|P|}\int_P|\nabla h|\,\dd x\le1+C\nu.
\]
Choose the final $\nu$ small enough that both errors are at most
$\eta_0$.  By \eqref{eq:repeat-slack}, the ratio of new selected
mass to old entire parent mass is then greater than
\[
 \frac{1+3\eta_0}{1+\eta_0}
 =1+\frac{2\eta_0}{1+\eta_0}.
\]
Fix $\zeta=\eta_0/(1+\eta_0)>0$.  This proves
\eqref{eq:local-mass-gain}; taking $\nu$ smaller still preserves
this value of $\zeta$.

\paragraph{Derivatives and the order of choices.}
Differentiating \eqref{eq:corrugation-formula} in the fixed
coordinates gives \eqref{eq:corrugation-derivatives}.  In a term
with $q$ derivatives on $\chi$, the scale factor is
\[
 s\ell(\nu S)^{-q}\ell^{-(j-q)}
 =s\ell^{1-j}\nu^q\le s\ell^{1-j}.
\]
All derivatives of the pattern and all powers of the bounded affine
matrix enter the fixed constant $C_j$.  The center slope and the
affine matrix are constants in the expression being differentiated.

To summarize the choices within this proof, first fix $F$ and its
small amplitude $a$, then the finite packing, cores, and parent
cutoff with the slack in \eqref{eq:repeat-slack}.  Choose the
preparation constants next, large and small respectively, and
finally choose a reciprocal $\nu$ satisfying the local positivity,
repeat preparation, and mass-error requirements.  This proves the
stated uniformity and completes the lemma.
\end{proof}

\subsection{Iteration and a common exponent for every finite order}

\begin{proposition}\label{prop:envelope}
After fixing the data of Lemma~\ref{lem:seed} and possibly shrinking
$\mathcal G$, there are common constants $\gamma>0$ and $c>0$
with the following property.  For every smooth $G\in\mathcal G$
and all sufficiently large $k$, there is a smooth real envelope
$h=h_k$ on $B$ such that $h=-a_*t+c_k$ on $\{t\le1/2\}$ and
\begin{equation}\label{eq:envelope-bounds}
 \begin{gathered}
  \|h\|_{C^j}\le C_j k^{\delta(j+1)}
  \quad\text{for every fixed integer }j\ge0,\\
  \max_{\substack{|e|_G=1\\\dd h(e)=0}}J_h(e)
       \ge c k^{-\delta}\quad\text{on }\{t\ge1/4\},
  \qquad
  \int_V\sqrt{1+|\nabla_Gh|_G^2}\,\dd x\ge c k^\gamma.
 \end{gathered}
\end{equation}
Here $\delta=10^{-6}$, and the chart $\Omega$ and region
$V\Subset\Omega$ are fixed.  All corrugations are supported in
one fixed initial cube compactly contained in $V$.
The constants $c_k$ may be any bounded prescribed real numbers.
The exponent $\gamma$ and all construction parameters are
independent of derivative orders, subsequent accuracy requirements,
and round gluing collars.  Constants $C_j$ may depend on the order,
the fixed smooth metric and coordinate systems, and the bound for
$|c_k|$; no uniform bound as $j\to\infty$ is asserted.
\end{proposition}

\begin{proof}
Choose $C_{\mathrm p}$ large enough and $c_{\mathrm p}>0$ small
enough for both Lemma~\ref{lem:corrugation} and the seed preparation
following Lemma~\ref{lem:seed}.  Choose the final $\nu$ sufficiently
small for the former lemma and so that a fixed cube $P_0$ of side
$\nu$ has closure in $U_0$.  All these choices are uniform for the
metrics in $\mathcal G$.  Put $r=\nu^2$.

Start with the prepared family $\mathcal P_0=\{P_0\}$ for $h_0$.
At generation $i\ge1$, apply Lemma~\ref{lem:corrugation} on every
cube of $\mathcal P_{i-1}$, using parent scale $S_i=r^{i-1}$,
and call the resulting function $h_i$.  The selected child cubes
form $\mathcal P_i$; their sides are $\nu r^i$ and they are
prepared at scale $r^i$.  Every generation consists of finitely
many cubes with disjoint interiors.  Each perturbation is supported
strictly inside its parent, so each $h_i$ is smooth, including across
all cube boundaries.  Nothing outside $P_0$ changes.

Let
\[
 I_i=\sum_{P\in\mathcal P_i}\int_P|\nabla h_i|\,\dd x.
\]
Equation~\eqref{eq:local-mass-gain} gives
\begin{equation}\label{eq:iterated-mass}
 I_i\ge(1+\zeta)I_{i-1},\qquad
 \int_V\sqrt{1+|\nabla h_i|^2}\,\dd x
 \ge I_i\ge I_0(1+\zeta)^i.
\end{equation}
The initial mass $I_0$ has a positive uniform lower bound by
Lemma~\ref{lem:seed} and the fixed volume of $P_0$.

We give the derivative accounting for all orders with one growth
constant.  Let $\beta\ge1$ and $0<\alpha\le1$ be the gradient
comparison factors of Lemma~\ref{lem:corrugation}.  The largest
gradient after $i$ generations is at most a fixed multiple of
$\beta^i$.  On $P_0$ every point has gradient at least
$s_{\min}\alpha^i$, since a point is changed at most once per generation.
At generation $i$, the center slope of each parent is therefore
at most $C\beta^{i-1}$.  By
\eqref{eq:corrugation-derivatives}, disjointness within that
generation gives
\begin{equation}\label{eq:generation-derivatives}
 \|D^j(h_i-h_{i-1})\|_{L^\infty}
 \le C_j\beta^{i-1}r^{i(1-j)}.
\end{equation}
There is no factor counting cubes in this supremum norm estimate.
There is also no product of higher pattern derivative constants
over generations: the increments are added, and their individual
formulas involve frozen slopes and affine matrices.  Passing to
the fixed finite atlas only adds fixed chain-rule constants and
lower-order derivatives, which satisfy the same upper bound.
Summation yields
\begin{align}
 \|h_m\|_{C^j}
 &\le C_j\left(1+\sum_{i=1}^m
                  \beta^{i-1}r^{i(1-j)}\right)\notag\\
 &\le C'_j\exp\left(
   [\log\beta+(j-1)_+|\log r|+1]m\right)
 \le C'_j e^{C_*(j+1)m},
 \label{eq:all-order-envelope-growth}
\end{align}
where $(j-1)_+=\max\{j-1,0\}$ and one finite $C_*$ works for
every $j\ge0$.

The lower Hessian margin has the same exponential control.
At a point outside $P_0$, the seed margin remains unchanged.  At a
point in $P_0$, consider the last generation $i\le m$ that changes
the function there.  If no such generation occurs, use the seed
margin.  Otherwise Lemma~\ref{lem:corrugation} bounds its retained
margin below by
\[
 c\,\frac{s_{\min}\alpha^{i-1}}{r^{i-1}}
 \ge c s_{\min}\alpha^m.
\]
Later perturbations vanish near that point, or all their jets there
vanish at their support boundaries, so this margin persists.  By
increasing $C_*$ if necessary,
\begin{equation}\label{eq:iterated-hessian-margin}
 \max_eJ_{h_m}(e)\ge c e^{-C_*m}
 \quad\text{on }\{t\ge1/4\}.
\end{equation}
This argument requires no nonzero-gradient assumption outside
$P_0$; critical points there are covered by the seed lemma.

Fix $\theta>0$ so small that $C_*\theta\le\delta$, and for each
$k\ge2$ take
\[
 m=\lfloor\theta\log k\rfloor,
 \qquad \gamma=\theta\log(1+\zeta)>0,
 \qquad h=h_m+c_k.
\]
Equations \eqref{eq:all-order-envelope-growth} and
\eqref{eq:iterated-hessian-margin} give the first two estimates in
\eqref{eq:envelope-bounds}.  Equation~\eqref{eq:iterated-mass}
gives the third, since
$(1+\zeta)^m\ge(1+\zeta)^{-1}k^\gamma$.
Adding bounded constants affects only the zeroth-order bound, and
$h=-a_*t+c_k$ on $t\le1/2$ because both the seed perturbation and
all corrugations lie above that set.

For clarity about subsequent uses, the constants determining
$C_*$, $\theta$, and $\gamma$ use only the fixed pattern, the
seed, and common $C^1$ metric bounds.  Higher derivatives of an
individual input metric may enter constants in a later smooth
coordinate change, but cannot alter these exponents.  Indeed, for
any fixed such coordinate change, an order-$j$ derivative of the
pullback is a finite sum of derivatives of $h$ of orders at most
$j$ multiplied by fixed smooth coefficients.  Its growth is still
bounded by $C_jk^{\delta(j+1)}$.  Each finite order may have its
own constant and frequency threshold.  This establishes the common
exponent for arbitrarily high finite accuracy and completes the
proof.
\end{proof}

\section{Local packets below an envelope}\label{sec:waves}

Throughout this Section the smooth positive pair \((A,\rho)\), and hence
\(G\) and \(L\), is fixed. The parameter \(k\) tends to infinity, and
\(h=h_k\) satisfies \eqref{eq:envelope-bounds}. Constants may depend on
the fixed pair, the constants in those bounds, and finitely many
construction orders. They are independent of \(k\) and of the packet
center. A lower threshold for \(k\) may depend on the same data.

The phase, transport amplitudes, and cutoff follow the complex WKB
construction of localized quasimodes; compare
\cite[Section~3]{DenckerSjostrandZworski2004}. Here the envelope varies
with the frequency. We therefore construct finite Taylor solutions and
track their dependence on its derivatives and Hessian margin, keeping
the support, derivative, and real-jet exponents fixed as the residual
accuracy increases.

We use complex bilinear extensions of real quadratic forms. When the
norm of a complex vector is used instead, it is the Hermitian norm.
In ordinary coordinate derivative estimates, a fixed finite atlas
with compactly contained subcharts is understood.

\subsection{The phase Hessian}

We first explain how the envelope test \(J_h\) enters the phase equation.
At a center \(z\), use \(G\)-normal coordinates and write
\(p=\nabla_Gh(z)\), \(Q=(1+|p|^2)^{1/2}\), and
\(H=\nabla_G^2h(z)\). For a unit vector \(v\perp p\), the complex
vector
\[
 P=p+iQv
\]
satisfies \(P\cdot P=-1\). A phase \(\phi\) with
\(\phi(z)=h(z)\) and \(\dd\phi(z)=P\) therefore satisfies the
constant term of the eikonal equation
\(G^{ij}\partial_i\phi\,\partial_j\phi=-1\). Its first-order terms
require \(SP=0\), where \(S=\partial^2\phi(z)\): the first metric
derivatives vanish in normal coordinates. A positive gap
\(H-\operatorname{Re}S\) then places \(\operatorname{Re}\phi\) below
\(h\) to quadratic order and supplies localization.

Write \(p=|p|e_1\) and \(b=|p|/Q\), choosing any unit
\(e_1\perp v\) when \(p=0\). For a symmetric matrix \(S=U+iD\),
the identity \(S(be_1+iv)=0\) forces
\[
 b^2U(e_1,e_1)+U(v,v)=0.
\]
Consequently, the corresponding test on the required gap is
\[
 b^2(H-U)(e_1,e_1)+(H-U)(v,v)
 =b^2H(e_1,e_1)+H(v,v)=J_h(v).
\]
This explains the positivity condition imposed on the envelope. The
following lemma constructs a compatible Hessian with a quantitative
gap, including when the real phase gradient vanishes.

\begin{lemma}[A dominated characteristic Hessian]\label{lem:phase-hessian}
Let \(e_1,v\) be orthonormal vectors in \(\R^4\), let \(0\le b<1\), and
let \(H\) be a real symmetric matrix such that
\[
 |H|\le k^{4\delta},\qquad
 J'=b^2H(e_1,e_1)+H(v,v)\ge k^{-2\delta}.
\]
For all sufficiently large \(k\), there are real symmetric matrices
\(U,D\) satisfying
\begin{equation}\label{eq:phase-hessian}
 (U+iD)(b e_1+i v)=0,\qquad
 H-U\ge k^{-3\delta}I,\qquad |U+iD|\le k^{20\delta}.
\end{equation}
The threshold can be chosen independently of \(H,e_1,v,b\) under these
hypotheses.
\end{lemma}

\begin{proof}
For \(b>0\), the first identity in \eqref{eq:phase-hessian} is equivalent
to
\[
 Dv=bUe_1,\qquad De_1=-Uv/b.
\]
The two prescribed columns are compatible with symmetry precisely when
\begin{equation}\label{eq:waves-column-compatibility}
 b^2U(e_1,e_1)+U(v,v)=0.
\end{equation}

First suppose that \(b\ge k^{-12\delta}\), and set
\[
 c_0=\frac{J'}{1+b^2},\qquad U=H-c_0I.
\]
Then \eqref{eq:waves-column-compatibility} holds. To give the extension
of \(D\) explicitly, write its blocks relative to
\(\R e_1\oplus\R v\oplus E\), where \(E=\{e_1,v\}^{\perp}\).
Denote the corresponding blocks of \(U\) by
\(u_{11},u_{12},u_{1E},u_{2E},\ldots\). Take
\[
 D=
 \begin{pmatrix}
 -u_{12}/b & b u_{11} & -u_{2E}/b\\
 b u_{11} & b u_{12} & b u_{1E}\\
 -u_{E2}/b & b u_{E1} & 0
 \end{pmatrix}.
\]
The identity \(u_{22}=-b^2u_{11}\) verifies both required columns.
Moreover,
\[
 H-U=c_0I\ge\tfrac12k^{-2\delta}I,\qquad
 |U|\le Ck^{4\delta},\qquad |D|\le Ck^{16\delta}.
\]
These imply \eqref{eq:phase-hessian} for large \(k\).

Now suppose that \(0\le b<k^{-12\delta}\). Put \(K=k^{14\delta}\) and
take
\[
 U=-K(I-vv^{\mathsf t})+b^2Kvv^{\mathsf t},\qquad
 D=-bK(e_1v^{\mathsf t}+ve_1^{\mathsf t}).
\]
The characteristic identity follows directly. These formulas also
give \(D=0\) and \(Uv=0\) when \(b=0\), so no division by a small \(b\)
is involved in this case.

For the domination inequality, split \(H-U-k^{-3\delta}I\) into
\(\R v\oplus v^\perp\). Its \(v^\perp\) block is at least
\(\tfrac12KI\) for large \(k\). Its scalar \(v\) block is at least
\[
 k^{-2\delta}-k^{-20\delta}-k^{-10\delta}-k^{-3\delta},
\]
because
\[
 H(v,v)\ge J'-b^2|H|,\qquad
 b^2|H|\le k^{-20\delta},\qquad b^2K\le k^{-10\delta}.
\]
The Schur complement subtracts at most
\(C|H|^2/K\le Ck^{-6\delta}\).
The positive \(k^{-2\delta}\) term dominates every subtraction.
The matrix is therefore positive definite for large \(k\).
Finally, \(|U|\le Ck^{14\delta}\) and \(|D|\le Ck^{2\delta}\),
proving the remaining bound.
\end{proof}

We will solve polynomial equations along a complex direction whose
bilinear square can be very small. Only its nonzero Hermitian norm
matters for this purpose.

\begin{lemma}[Polynomial directional inverses]\label{lem:polynomial-inverse}
Let \(\mathcal P_m\) be the space of complex homogeneous polynomials of
degree \(m\) on \(\R^4\), equipped with the coefficient norm in a fixed
monomial basis. Fix \(0<c\le C<\infty\).
If \(q\in\C^4\) satisfies \(c\le |q|\le C\), then
\[
 q\cdot\partial:\mathcal P_{m+1}\longrightarrow\mathcal P_m
\]
has a right inverse with norm bounded in terms of \(m,c,C\).
No lower bound for \(|q\cdot q|\) is required.
\end{lemma}

\begin{proof}
Choose an index \(i\) with \(|q_i|\ge c/2\), write
\(q\cdot\partial=q_i\partial_i+D'\), and let \(I_i\) denote integration
in \(X_i\) with zero integration constant. On \(\mathcal P_m\), set
\[
 \mathcal K=\frac{D'I_i}{q_i},\qquad
 R_{q,m}=\frac{I_i}{q_i}\sum_{r=0}^{m}(-\mathcal K)^r.
\]
Each application of \(\mathcal K\) lowers the degree in the variables
other than \(X_i\). Thus \(\mathcal K^{m+1}=0\) on \(\mathcal P_m\).
Since
\[
 (q\cdot\partial)\frac{I_i}{q_i}=I+\mathcal K,
\]
the finite geometric sum gives \((q\cdot\partial)R_{q,m}=I\).
The stated coefficient bound follows from the finite formula and the
bounds on \(q_i^{-1}\) and \(q\). The four possible choices of \(i\)
have the same uniform bound.
\end{proof}

\subsection{Construction and estimates for a packet family}

For a complex packet \(w\) centered at \(z\), let
\[
 \mathbf j_z w(x)
 =e^{-kh(x)}
   \left(w(x),\frac{\partial_\xi w(x)}{kQ_z}\right)\in\C^5,\qquad
 Q_z=\sqrt{1+|\nabla h(z)|^2}.
\]
Here \(\xi\) denotes the single normal coordinate system chosen at \(z\);
the denominator \(Q_z\) is kept fixed when \(x\) varies. A family
\(\mathcal W_z\) determines a real \(5\)-by-\(2|\mathcal W_z|\) matrix
\(\mathcal J_z(x)\) whose columns are the real and imaginary parts of
these jets. Its least spanning singular value is the square root of
the least eigenvalue of
\(\mathcal J_z(x)\mathcal J_z(x)^{\mathsf t}\).

\begin{proposition}[Packets with arbitrary finite accuracy]\label{prop:packets}
Fix a nonnegative integer \(J\) and a real number \(M>0\). Finite
construction orders can be chosen, before taking \(k\) large, so that
every center \(z\) with \(t(z)\ge1/4\) has a family \(\mathcal W_z\) of
four smooth complex packets with the following properties.

Each packet is supported in an ordinary coordinate ball about \(z\)
of radius at most \(Ck^{-1/3-100\delta}\). For every fixed
nonnegative integer \(j\),
\begin{equation}\label{eq:packet-derivatives}
 |D_x^j w(x)|\le C_j k^{2j}e^{kh(x)}.
\end{equation}
For the prescribed orders,
\begin{equation}\label{eq:packet-residual}
 |D_x^j(L+k^2)w(x)|\le C e^{kh(x)}k^{-M},
 \qquad 0\le j\le J.
\end{equation}
The center jets and the nearby jets satisfy
\begin{equation}\label{eq:jet-span}
 \sigma_{\min}(\mathcal J_z(z))\ge c k^{-3/4},\qquad
 \sigma_{\min}(\mathcal J_z(x))\ge k^{-1}
 \quad\text{if }\operatorname{dist}(x,z)\le k^{-10},
\end{equation}
for sufficiently large \(k\). Any fixed smooth reference distance
may be used here. The same conclusion holds if the latter radius
is multiplied by a fixed constant.

The exponents in the support, derivative, and jet bounds are
independent of \(J,M\). Constants and lower frequency thresholds may
depend on the chosen finite orders.
\end{proposition}

\begin{proof}
We divide the construction into the selection of jets, finite
polynomial equations, and localization estimates.

\paragraph{Normal coordinates and the linear jets.}
For each center choose a \(G\)-orthonormal frame and use normal
coordinates \(\xi\) in that frame. Put
\[
 p=\nabla h(z),\qquad Q=Q_z,\qquad H=\nabla^2h(z).
\]
In these coordinates the tangent and cotangent components agree at
the center. The envelope bounds, with a little slack in their
exponents, give
\begin{equation}\label{eq:waves-center-bounds}
 Q\le k^{3\delta},\qquad |H|\le k^{4\delta}
\end{equation}
for large \(k\).

If \(p\ne0\), set \(e_1=p/|p|\), and choose a favorable unit test \(e_2\)
from \eqref{eq:envelope-bounds}. If \(p=0\), first choose such an
\(e_2\), then choose any unit \(e_1\perp e_2\).
Complete to an orthonormal basis \(e_1,e_2,e_3,e_4\), and put
\[
 \eta_k=k^{-9\delta},\qquad
 v_2=e_2,\quad
 v_3=\frac{e_2+\eta_k e_3}{\sqrt{1+\eta_k^2}},\quad
 v_4=\frac{e_2+\eta_k e_4}{\sqrt{1+\eta_k^2}}.
\]
These three vectors span \(e_1^\perp\) with least spanning singular
value at least \(c\eta_k\).

Use the four pairs
\[
 (p',v)=(p,v_2),(p,v_3),(p,v_4),
 \bigl(p+Qk^{-3/4}e_1,v_2\bigr).
\]
For each pair define
\[
 Q'=\sqrt{1+|p'|^2},\qquad P=p'+iQ'v,\qquad b=|p'|/Q'.
\]
Then \(p'=|p'|e_1\), \(P\cdot P=-1\), and
\begin{equation}\label{eq:waves-imaginary-size}
 1\le Q'/Q\le1+k^{-3/4}.
\end{equation}
Changing \(e_2\) to a \(v_j\) changes its Hessian value by at most
\(Ck^{-5\delta}\). If \(a=|p|\), the function \(a^2/(1+a^2)\)
changes by at most \(2k^{-3/4}\) when \(a\) is increased by
\(Qk^{-3/4}\). Indeed its derivative at \(s\ge a\) is at most
\(2(1+a^2)^{-3/2}\). The additional error in the Hessian test is
at most \(Ck^{-3/4+4\delta}\). Both errors are negligible compared
with the original \(ck^{-\delta}\) margin. Consequently
\[
 b^2H(e_1,e_1)+H(v,v)\ge k^{-2\delta}
\]
for every selected pair. Apply Lemma~\ref{lem:phase-hessian} to obtain
\(U,D\). Multiplying its characteristic identity by \(Q'\) gives
\((U+iD)P=0\).

\paragraph{Frozen charts and scaling.}
The choice of frame may depend on \(k\), but each packet uses one
fixed chart when spatial derivatives are taken. More explicitly, if
\(E\) is its chosen frame, the chart map is
\[
 \Psi_{z,E}(\xi)=\exp_z\left(\sum_i\xi_iE_i\right).
\]
For each fixed derivative order \(q\), the derivatives of these maps
and their inverses are uniformly bounded over all centers and
orthonormal frames on sufficiently small common coordinate balls.
This follows from smoothness and compactness of the orthonormal
frame bundle of the fixed metric. In particular, the chain rule
gives
\[
 |D_\xi^q(h\circ\Psi_{z,E})|
 \le C_q(G)\sum_{\ell=1}^q\|h\|_{C^\ell}
 \le C'_q k^{\delta(q+1)}.
\]
There are no spatial derivatives of the rule used to select \(E\).

Set
\begin{equation}\label{eq:waves-scaling}
 r_0=k^{-100\delta},\qquad X=\xi/r_0,\qquad
 \kappa=kr_0Q,\qquad
 h_{\rm sc}(X)=\frac{h(\Psi_{z,E}(r_0X))-h(z)}{r_0Q}.
\end{equation}
For every fixed \(q\ge2\),
\begin{equation}\label{eq:waves-scaled-envelope}
 |D_X^q h_{\rm sc}|
 \le C_q Q^{-1}k^{\delta(101-99q)}.
\end{equation}
All these derivatives are bounded independently of \(k,z\). The
constant may depend on \(q\) and on high derivatives of the fixed
metric.

In normal coordinates write
\(L=g^{ij}(\xi)\partial_{\xi_i\xi_j}
+\beta^i(\xi)\partial_{\xi_i}\).
After multiplication by \(r_0^2\), the operator \(L+k^2\) becomes
\[
 \mathcal L+\kappa^2/Q^2,\qquad
 \mathcal L=g^{ij}(r_0X)\partial_{X_iX_j}
        +r_0\beta^i(r_0X)\partial_{X_i}.
\]
Its coefficients have uniformly bounded derivatives of every fixed
order on a fixed \(X\)-ball. The normal-coordinate identities are
\(g^{ij}(0)=\delta_{ij}\) and
\(\partial_{\xi_\ell}g^{ij}(0)=0\).

\paragraph{Finite eikonal and transport equations.}
Choose an integer \(T\ge3\), to be specified below. We construct a
complex polynomial \(\Phi\) with
\[
 \Phi(0)=0,\qquad
 \partial_X\Phi(0)=P/Q,\qquad
 \partial_X^2\Phi(0)=r_0(U+iD)/Q
\]
such that
\[
 E_\Phi(X)
 =g^{ij}(r_0X)\partial_i\Phi\,\partial_j\Phi+Q^{-2}
\]
vanishes to order \(T\) at zero. Here \emph{vanishes to order \(T\)}
means that all derivatives of order less than \(T\) vanish.
The constant equation follows from \(P\cdot P=-1\), and the linear
equations follow from \((U+iD)P=0\) and the vanishing first derivatives
of the metric.

At homogeneous degree \(m\ge2\), the as-yet unspecified homogeneous
term \(\Phi_{m+1}\) enters the equation through
\(2(P/Q)\cdot\partial\,\Phi_{m+1}\). All other terms have already
been determined. The Hermitian norm of \(P/Q\) lies between \(1\)
and \(3\) for large \(k\), by \eqref{eq:waves-imaginary-size}.
Lemma~\ref{lem:polynomial-inverse} therefore solves each equation
with a uniformly bounded right inverse. Solving through degree
\(T-1\) gives the desired polynomial. Its coefficients are bounded
for each fixed \(T\): the prescribed Hessian is bounded by
\(k^{-80\delta}\), the coefficient jets of \(\mathcal L\) are bounded,
and only finitely many bounded algebraic operations are used.

Define
\[
 \mathcal T a
 =2g^{ij}(r_0X)(\partial_i\Phi)(\partial_ja)
       +(\mathcal L\Phi)a.
\]
Fix an integer \(s\ge0\). Choose polynomials \(a_0,\ldots,a_s\) such
that
\[
 a_0(0)=1,\qquad a_j(0)=0\quad(j\ge1),
\]
and each of
\[
 \mathcal Ta_0,\qquad
 \mathcal Ta_j+\mathcal La_{j-1}\quad(1\le j\le s)
\]
vanishes to order \(T\) at zero. To see this directly, solve
\(\mathcal Ta=f\) by homogeneous degrees. The new amplitude term
of degree \(m+1\) enters the degree-\(m\) equation through
\(2(P/Q)\cdot\partial\); the remaining terms use only coefficients
already chosen. Lemma~\ref{lem:polynomial-inverse} again applies.
First construct the entire \(a_0\), then \(a_1\), and so on.
Thus the second derivative in \(\mathcal La_{j-1}\) causes no circular
dependence. Constants at zero are free, as stipulated above.
For fixed \(T,s\), all polynomial coefficients are bounded uniformly
in \(k,z\). No bound uniform as \(T,s\to\infty\) is asserted or needed.

Choose once a real smooth cutoff \(\omega\), equal to one on
\(\{|X|\le1\}\), supported in \(\{|X|<2\}\), and taking values in
\([0,1]\). Define
\begin{equation}\label{eq:packet-definition}
 \begin{split}
 a(X)&=\sum_{j=0}^s\kappa^{-j}a_j(X),\\
 w(\Psi_{z,E}(r_0X))
 &=e^{kh(z)}e^{\kappa\Phi(X)}a(X)\omega(k^{1/3}X),
 \end{split}
\end{equation}
and extend \(w\) by zero. For large \(k\), its support is compactly
contained in the normal chart, so this extension is smooth and has
the asserted support radius. On \(|X|\le2k^{-1/3}\),
\begin{equation}\label{eq:waves-amplitude}
 a(X)=1+O(k^{-1/3}),\qquad |a(X)|\ge\tfrac12,
\end{equation}
where the constant depends on \(T,s\). This follows from the chosen
constant terms and \(\kappa\ge1\).

\paragraph{Localization.}
The linear part of
\(\operatorname{Re}\Phi-h_{\rm sc}\) is either zero or
\(k^{-3/4}e_1\cdot X\). Lemma~\ref{lem:phase-hessian},
\eqref{eq:waves-center-bounds}, and
\eqref{eq:waves-scaled-envelope} imply
\begin{equation}\label{eq:waves-localization}
 \operatorname{Re}\Phi(X)-h_{\rm sc}(X)
 \le k^{-3/4}|X|-c k^{-106\delta}|X|^2+C|X|^3.
\end{equation}
The third derivative bounds used here are uniform for each fixed
construction order. On the packet support the cubic term can be
absorbed into half the quadratic deficit, since
\(k^{-1/3+106\delta}\to0\).
The maximum of the remaining linear term minus the quadratic term
is \(O(k^{-3/2+106\delta})\). Also
\[
 k^{1-100\delta}\le\kappa\le k.
\]
Consequently
\[
 |e^{kh(z)+\kappa\Phi(X)}|
 \le 2e^{kh(\Psi_{z,E}(r_0X))}
 \quad (|X|\le2k^{-1/3})
\]
for large \(k\).

On the smaller annulus
\(\tfrac14k^{-1/3}\le|X|\le2k^{-1/3}\), the negative exponent has
magnitude at least \(ck^{1/3-206\delta}\), whereas the positive linear
exponent is at most \(Ck^{-1/12}\). Since
\(1/3-206\delta>1/4\), a larger threshold gives
\begin{equation}\label{eq:waves-exponential-tail}
 |e^{kh(z)+\kappa\Phi(X)}|
 \le e^{-2k^{1/4}}e^{kh(\Psi_{z,E}(r_0X))}
 \quad\left(\tfrac14k^{-1/3}\le|X|\le2k^{-1/3}\right).
\end{equation}

\paragraph{Residual and derivative estimates.}
For the uncut expression, direct differentiation gives
\[
 \begin{split}
 e^{-\kappa\Phi}
 (\mathcal L+\kappa^2/Q^2)(e^{\kappa\Phi}a)
 ={}&\kappa^2E_\Phi a+\kappa\mathcal Ta_0\\
 &+\sum_{j=1}^s\kappa^{1-j}
       (\mathcal Ta_j+\mathcal La_{j-1})
       +\kappa^{-s}\mathcal La_s.
 \end{split}
\]
Taylor's theorem bounds a remainder vanishing to order \(T\), after
\(q\le J\le T\) derivatives in \(X\), by \(C|X|^{T-q}\).
On the support this is at most \(Ck^{-(T-q)/3}\).
Each ordinary derivative of such a remainder therefore has a cost
at most \(k^{1/3+100\delta}\), relative to its zeroth-order Taylor
bound. A derivative falling on the exponential, amplitude, or cutoff
has a cost at most a fixed constant times
\[
 r_0^{-1}(1+\kappa+k^{1/3})\le Ck^{1+3\delta}<Ck^2.
\]
Higher coordinate derivatives obey the same conservative \(k^2\)
cost per derivative, by the chain rule and the fixed normal-chart
bounds.

Undoing the operator scaling costs \(r_0^{-2}=k^{200\delta}\).
Thus, for \(0\le j\le J\), the uncut residual divided by \(e^{kh}\)
is bounded by
\begin{equation}\label{eq:waves-accuracy-budget}
 C k^{200\delta+2J}
 \left(k^{2-T/3}+k^{1-T/3}
       +k^{-(1-100\delta)s}\right).
\end{equation}
The commutators with the cutoff, and any prescribed derivatives of
them, have polynomial factors times \eqref{eq:waves-exponential-tail}.
They are smaller than every inverse power of \(k\).

For example, it suffices to choose integers \(T,s\) with
\[
 T>3(M+200\delta+2J+3),\qquad
 s>\frac{M+200\delta+2J+1}{1-100\delta}.
\]
These choices are made before the frequency threshold.
Increasing that threshold absorbs all their finite constants and
proves \eqref{eq:packet-residual}. The same derivative calculation
without the operator factor proves \eqref{eq:packet-derivatives} for
every fixed \(j\), with exponent \(2j\) independent of \(T,s\).

\paragraph{Real jet spanning.}
At the center, \(a(0)=1\) and the cutoff is constant, so
\[
 \mathbf j_z w(z)=\left(1,\frac{P}{Q}\right)+O(\kappa^{-1}).
\]
The error is bounded by \(Ck^{-1+100\delta}\).
Before this error, the real parts contain the two vectors
\[
 (1,(|p|/Q)e_1),\qquad
 (1,(|p|/Q+k^{-3/4})e_1).
\]
Their coefficients are bounded, and their least spanning singular
value on the value--\(e_1\) plane is at least \(ck^{-3/4}\).
The imaginary parts span the orthogonal \(e_1^\perp\) slope space
with least singular value at least \(ck^{-9\delta}\).
Together their least spanning singular value is at least
\(ck^{-3/4}\). Since
\(k^{-1+100\delta}=o(k^{-3/4})\), the actual center jets retain this
bound, with a smaller constant. The argument applies without
change at \(p=0\).

Finally, \eqref{eq:packet-derivatives} through order two and the
first envelope derivative bound give
\[
 |D_x(\mathbf j_z w(x))|\le Ck^3
\]
near the center. In differentiating this expression \(Q_z\) remains
fixed, while differentiating \(e^{-kh(x)}\) costs
\(k|D_xh|\le Ck^{1+2\delta}\).
A displacement of \(O(k^{-10})\) therefore changes the jet matrix
by \(O(k^{-7})\). Uniformly bounded changes between ordinary and
normal coordinates preserve these estimates.
This is negligible relative to \(k^{-3/4}\), and proves the second
inequality of \eqref{eq:jet-span}, after increasing the threshold.
\end{proof}

\subsection{Deterministic estimates at the local oscillation scale}

The next Lemma supplies estimates for an arbitrary polynomial-sized
collection of the preceding packets. It includes packets whose
centers lie outside the region in which signs will be measured.
Use the fixed chart \(\Omega\subset\{t>1/2\}\) and the open set
\(V\Subset\Omega\) from the preceding Sections. Points of \(\Omega\)
are identified with their coordinate vectors.

\begin{lemma}[Packets on a wavelength ball]\label{lem:scaled-packets}
Fix the finite construction orders in Proposition~\ref{prop:packets}
and a fixed finite number \(K_1\ge0\). Let \(\mathcal W\) be the collection of
all four packets from each of finitely many centers in
\(\{t\ge1/4\}\), with \(|\mathcal W|\le Ck^{K_1}\).
Suppose that for every \(x\in V\) some center is within
\(C_0k^{-10}\) of \(x\), where \(C_0\) is fixed. Set
\[
 R_x=(kQ_x)^{-1},\qquad
 S_x^2=\sum_{w\in\mathcal W}|w(x)|^2,\qquad
 w_x(Y)=w(x+R_xY),\quad |Y|\le2.
\]
The coordinate ball in this definition lies in \(\Omega\) for large
\(k\).

Consider the packets whose supports meet this closed ball. Their
centered scaled coordinates \(X_z\) are defined on the whole ball.
Partition them into
\[
 \mathcal N_x=\{w:|X_z(x)|<\tfrac12k^{-1/3}\},\qquad
 \mathcal T_x=\{w:|X_z(x)|\ge\tfrac12k^{-1/3}\}.
\]
For every fixed nonnegative integer \(j\), uniformly in \(x\in V\),
\begin{equation}\label{eq:scaled-packet-bounds}
 \begin{aligned}
 S_x&\ge c e^{kh(x)},\\
 \sup_{|Y|\le2}|D_Y^j w_x(Y)|
     &\le C_j|w(x)| &&(w\in\mathcal N_x),\\
 \sup_{|Y|\le2}|D_Y^j w_x(Y)|
     &\le C_j e^{-k^{1/4}}e^{kh(x)}
                                  &&(w\in\mathcal T_x).
 \end{aligned}
\end{equation}
Every packet in \(\mathcal N_x\) has cutoff identically one and
nonzero amplitude on the ball. Moreover,
\begin{equation}\label{eq:imaginary-gradient}
 \left|\operatorname{Im}
   \left(\frac{\dd_Yw_x}{w_x}\right)(0)\right|\ge c
 \qquad(w\in\mathcal N_x).
\end{equation}
In particular,
\begin{equation}\label{eq:waves-square-sums}
 \begin{aligned}
 \frac1{S_x^2}\sum_{w\in\mathcal W}
   \sup_{|Y|\le2}|D_Y^jw_x(Y)|^2&\le C_j,\\
 \frac1{S_x^2}\sum_{w\in\mathcal N_x}|w(x)|^2&=1-o(1).
 \end{aligned}
\end{equation}
The constants and thresholds may depend on the fixed construction
orders and on \(K_1,C_0\), but not on \(k,x\).
\end{lemma}

\begin{proof}
The bounds on the first two derivatives of \(h\) imply
\[
 |D_xQ|\le Ck^{4\delta}
\]
in the fixed chart. A packet centered at \(z\) that meets the ball
about \(x\) has
\[
 |z-x|\le Ck^{-1/3-100\delta}+C/k.
\]
For large \(k\), such a \(z\), and the whole ball, belong to a
compactly contained part of \(\Omega\). Normal coordinates centered
at \(z\) are defined throughout the ball. The preceding estimates
give
\[
 |Q_z-Q_x|
 \le Ck^{-1/3-96\delta}+Ck^{-1+4\delta}=o(1).
\]
In particular, \(Q_z/Q_x\) is bounded above and below independently
of the packet. The same comparison holds for \(Q\) throughout the
ball. Since \(|\dd h|\le CQ\) in the coordinate chart,
\begin{equation}\label{eq:waves-local-weight}
 k|h(x+R_xY)-h(x)|\le C,\qquad |Y|\le2.
\end{equation}

The nearest center has
\(|X_z(x)|=O(k^{-10+100\delta})\).
The center formula and its nearby-jet estimate in the proof of
Proposition~\ref{prop:packets} show that one of its packets has
\(|w(x)|\ge c e^{kh(x)}\). This proves the first line of
\eqref{eq:scaled-packet-bounds}.

Write \(\kappa_z=kr_0Q_z\). For every fixed \(q\ge1\), the chain
rule and normal-chart bounds give
\[
 |D_Y^qX_z|\le C_qr_0^{-1}R_x^q,\qquad
 \kappa_z|D_Y^qX_z|
 \le C_q\,\frac{Q_z}{Q_x}(kQ_x)^{1-q}\le C_q.
\]
For \(q=1\) there is also a lower bound: the matrix
\begin{equation}\label{eq:waves-scaled-jacobian}
 \kappa_z D_YX_z=(Q_z/Q_x)D_x\xi_z
\end{equation}
has inverse bounded independently of \(k,x,z\).
In particular,
\[
 |X_z(x+R_xY)-X_z(x)|
 \le Ck^{-1+100\delta}/Q_x=o(k^{-1/3}).
\]

Suppose first that \(w\in\mathcal T_x\).
At every point where its cutoff or any derivative of the cutoff can
contribute,
\[
 \tfrac14k^{-1/3}\le |X_z|\le2k^{-1/3}.
\]
Apply \eqref{eq:waves-exponential-tail}. Each fixed positive
\(Y\)-derivative of \(\kappa_z\Phi(X_z)\) is bounded: in a term
containing \(m\) derivatives of \(X_z\), the displayed derivative
bounds leave a factor \(\kappa_z^{1-m}\le1\).
Amplitude derivatives are bounded as well.
Derivatives of \(\omega(k^{1/3}X_z)\) are bounded because
\(k^{1/3}/\kappa_z\le k^{-2/3+100\delta}\le1\).
Thus differentiation introduces only fixed constants at each fixed
order. Together with \eqref{eq:waves-local-weight}, this proves
the tail bound in \eqref{eq:scaled-packet-bounds}.
This argument uses an absolute estimate and is valid even if
\(w(x)=0\).

For \(w\in\mathcal N_x\), the displacement estimate shows that
\(|X_z|<k^{-1/3}\) throughout the ball, so its cutoff is identically
one. By \eqref{eq:waves-amplitude}, the amplitude is bounded away
from zero. The bounds just proved for derivatives of
\(\kappa_z\Phi(X_z)\) and \(a(X_z)\) imply bounded derivatives of
\(w_x\) relative to \(|w_x|\).
Integrating the bounded real logarithmic derivative along straight
segments in the \(Y\)-ball gives
\[
 |w_x(Y)|\le C|w_x(0)|=C|w(x)|.
\]
This proves the second line of \eqref{eq:scaled-packet-bounds},
including \(j=0\).

For the lower oscillation estimate, the bounded coefficients of
\(\Phi\) give
\[
 \operatorname{Im}\partial_X\Phi(X_z(x))
   =(Q'/Q_z)v+O(k^{-1/3}).
\]
By \eqref{eq:waves-imaginary-size} its leading term has norm at
least one. The inverse bound in \eqref{eq:waves-scaled-jacobian}
therefore gives a positive lower bound for the imaginary
\(Y\)-gradient of \(\kappa_z\Phi(X_z)\).
The amplitude contributes only
\[
 \left|\frac{\dd_Ya(X_z)}{a(X_z)}\right|
 \le C r_0^{-1}R_x=Ck^{-1+100\delta}/Q_x=o(1).
\]
It cannot cancel that lower bound.
This proves \eqref{eq:imaginary-gradient}, including at centers
with \(p=0\) and for phases nearly isotropic in the complex
bilinear form.

Packets whose supports do not meet the ball contribute zero to
all the displayed sums. For the other packets, squaring
\eqref{eq:scaled-packet-bounds}, dividing by
\(S_x^2\ge c e^{2kh(x)}\), and using the cardinality bound shows
that the total tail contribution is at most
\[
 C_j k^{K_1}e^{-2k^{1/4}}=o(1).
\]
The non-tail contributions are bounded by a constant times
\(\sum_{\mathcal N_x}|w(x)|^2/S_x^2\le1\).
At derivative order zero the definition of \(S_x\) also shows that
this latter ratio equals \(1-o(1)\).
This proves \eqref{eq:waves-square-sums}.
\end{proof}

\section{Exact eigenfunctions and a fixed metric}\label{sec:global}

We combine the packets with the round beam and choose a real sample
with a nonvanishing first jet and a large sign certificate.
Quantitative control of the first jet permits exact correction of the
coefficients; the certificate preserves the nodal lower bound under later
perturbations. We then make the eigenvalue simple in the remaining
round collar and use spectral persistence to pass to one smooth pair.

We use the common neighborhood, regions $V\Subset\Omega$, and exponent
$\gamma>0$ supplied by Proposition~\ref{prop:envelope}. Shrink the
neighborhood of positive pairs, if necessary, so that its associated
metrics belong to the neighborhood required there. Throughout a single
step the input pair is fixed. Constants may depend on this pair, on a
fixed positive collar width, and on fixed finite approximation orders.
Their frequency exponents, when specified below, do not have these
dependencies.

\subsection{A weighted Gaussian construction}

Fix $l_-,l_+$ as in Proposition~\ref{prop:step}, put $w=l_+-l_->0$,
and choose
\[
 s_i=l_-+iw/5\quad(1\le i\le4).
\]
For the allowed frequency $k=\sqrt{N(N+3)}$, define
\[
 b_N(t)=\frac{N}{2k}\log(1-t),\qquad t<1.
\]
Apply Proposition~\ref{prop:envelope} and choose the additive constant
in $h$ so that $h(s_2)=b_N(s_2)$. Here the notation $h(s_2)$ is
unambiguous because $h=-a_*t+\text{constant}$ on $t\le1/2$.
The constants added in this way are bounded, since
$s_2\in[1/4,1/2]$ and $N/k\le1$. For the fixed sufficiently small
seed slope $a_*$ and all sufficiently large $N$,
\[
 \frac{\dd}{\dd t}(h-b_N)
 =-a_*+\frac{N}{2k(1-t)}\ge c_0>0
 \qquad(0\le t\le1/2),
\]
where $c_0$ can be fixed independently of the collar. Consequently,
\begin{equation}\label{eq:gluing-gaps}
 h\le b_N-cw\quad(t\le s_1),\qquad
 h\ge b_N+cw\quad(s_3\le t\le s_4)
\end{equation}
for a fixed $c>0$.

Choose a smooth cutoff $0\le\beta\le1$ equal to one for $t\le s_3$
and zero for $t\ge s_4$. Set
\begin{equation}\label{eq:gluing-weight}
 W=e^{kh}+\beta(t)e^{kb_N(t)}.
\end{equation}
The second summand is extended by zero outside $t<s_4$; this is
smooth. The input pair is round on the support of $\beta$. Thus the
only residual of $\beta u_*$ is produced by derivatives of $\beta$,
where the second inequality in \eqref{eq:gluing-gaps} applies. For
every fixed nonnegative integer $q$ and every fixed $M>0$,
\begin{equation}\label{eq:beam-weighted-bounds}
 \begin{split}
 |D^q(\beta u_*)|&\le C_q k^qW,\\
 |D^q(L+k^2)(\beta u_*)|&\le C_{q,M}k^{-M}W,\\
 |\dd W|/W&\le Ck^2 .
 \end{split}
\end{equation}
All derivatives here and below are in the fixed finite coordinate
systems. To justify the first bound, work on $t\le s_4<1/2$ and use
$u_*=(1-t)^{N/2}\cos(N\vartheta)$ in finitely many angular charts.
Each derivative costs at most a fixed constant times $k$. In the
cutoff region the exponential gap absorbs the cutoff derivatives.
The second bound follows in the same way from the commutator with
$\beta$. Finally, $|\dd h|\le Ck^{2\delta}$ and the same gap control
the derivative of \eqref{eq:gluing-weight}, giving its last bound.
For a fixed collar, all derivatives of $\beta$ are fixed constants;
more quantitatively they may be bounded by $C_qw^{-q}$.

Choose a net in $\{t\ge s_1\}$ of mesh at most a fixed small multiple
of $k^{-10}$, with at most $Ck^{40}$ centers. At every center take the
fixed finite family of Proposition~\ref{prop:packets}. Denote the
resulting collection, counting packets with their center and family
index, by $\mathcal W_k$. Its cardinality is at most $Ck^{40}$.
For each $v\in\mathcal W_k$ take two independent standard real
Gaussian variables $g_v^{(1)},g_v^{(2)}$, all independent, and set
\begin{equation}\label{eq:gaussian-sum}
 u=\beta u_*+
 \sum_{v\in\mathcal W_k}
 \bigl(g_v^{(1)}\operatorname{Re}v+
       g_v^{(2)}\operatorname{Im}v\bigr).
\end{equation}
Let $\mathcal B_k$ be the event that all these coefficients have
absolute value at most $k$. A Gaussian tail estimate gives
\begin{equation}\label{eq:coefficient-event}
 \Pp(\mathcal B_k^c)\le Ck^{40}e^{-k^2/2}\longrightarrow0.
\end{equation}

Given a finite integer $J\ge2$ and a number $M>0$, choose the packet
orders so that their residuals through order $J$ decay at least as
$k^{-M-42}e^{kh}$. This choice is made before $k$ is taken large.
The packet estimates, \eqref{eq:beam-weighted-bounds}, and the
cardinality bound imply, on $\mathcal B_k$,
\begin{equation}\label{eq:sum-weighted-bounds}
 \begin{split}
 |D^qu|&\le C_q k^{2q+41}W\qquad(0\le q\le J),\\
 |D^q\mathfrak r|&\le C k^{-M}W\qquad(0\le q\le J),\\
 \mathfrak r&=\diver_\mu(A\,\dd u)+k^2\rho u .
 \end{split}
\end{equation}
The multiplication by the fixed smooth $\rho$ is included in the
second estimate. In particular, increasing $J$ or $M$ does not change
the exponent $2q+41$ in the first estimate.

Every packet is supported within distance $Ck^{-1/3-100\delta}$ of
its center. The function $t$ has bounded differential, so for a fixed
$C$ all packets vanish on
\begin{equation}\label{eq:untouched-region}
 t<s_1-Ck^{-1/3-100\delta}.
\end{equation}
On this set $u=u_*$ and $\mathfrak r=0$.

We use the full five-dimensional jet spanning in \eqref{eq:jet-span}
to show that, with probability tending to one, $u$ and its first
derivatives never become too small simultaneously. A lower bound by
an inverse power of $k$ is sufficient:
the coefficient correction will spend a larger power of residual
accuracy to compensate for it.

\begin{lemma}[A quantitative nonvanishing first jet]
\label{lem:noncancellation}
For every fixed choice of the finite packet orders, the construction
above satisfies
\begin{equation}\label{eq:noncancellation}
 \Pp\left(\mathcal B_k\cap\left\{
 \inf_{x\in B}\frac{|u(x)|+|\dd u(x)|_{g_0}}{W(x)}
 \ge k^{-500}\right\}\right)\longrightarrow1.
\end{equation}
The exponent $500$ is independent of those orders and of the fixed
input pair and collar.
\end{lemma}

\begin{proof}
We separate two deterministic regions. If $W/e^{kh}>k^{50}$, then
$\beta=1$ for large $k$: on its transition region
\eqref{eq:gluing-gaps} instead gives $W/e^{kh}\le1+e^{-cw k}$.
In the uncut region, the cosine formula and the angular derivative
give
\[
 |u_*|+|\dd u_*|_{g_0}\ge c e^{kb_N}.
\]
On $\mathcal B_k$, the packet value and first derivative together are
bounded by $Ck^{43}e^{kh}$. The beam therefore dominates their sum
on this region and gives a lower bound $cW$ for the first jet.
The first inequality in \eqref{eq:gluing-gaps} puts all of $t<s_1$
in this region for sufficiently large $k$.

At a point in the remaining region there is a nearby net center.
The persisting real jet span in \eqref{eq:jet-span}, after converting
the normalized derivative to ordinary coordinate derivatives, gives
for the value and four derivatives of the random sum a covariance
matrix whose least eigenvalue has square root at least
\[
 c k^{-1}e^{kh}\ge c k^{-51}W\ge k^{-52}W
\]
for all sufficiently large $k$. Only one nearby packet family is
needed for this lower bound. All other families add nonnegative
covariance matrices. The deterministic beam is a mean shift and
does not affect the maximum of a Gaussian density. Hence, in any
fixed coordinate chart and for any fixed constant $C_0$, the
\emph{unconditional} probability that the normalized jet
\[
 \mathcal J(x)=W(x)^{-1}(u(x),\partial_1u(x),\ldots,\partial_4u(x))
\]
has norm at most $C_0k^{-D}$ is bounded by
\begin{equation}\label{eq:jet-small-ball}
 C k^{260-5D}.
\end{equation}
Indeed the density is at most $Ck^{260}$ and the five-dimensional
ball has volume $Ck^{-5D}$. No conditioning on $\mathcal B_k$ is used
in this estimate.

On $\mathcal B_k$, \eqref{eq:sum-weighted-bounds} through order two
and the last bound in \eqref{eq:beam-weighted-bounds} give
\[
 |D\mathcal J|\le Ck^{45}\le k^{46}
\]
for sufficiently large $k$. In each of finitely many coordinate
patches covering $B$, take a grid of mesh $k^{-D-48}$. Its total
cardinality is at most $Ck^{4D+192}$. Apply
\eqref{eq:jet-small-ball} at grid points in the Gaussian region;
on the other grid points the deterministic beam bound already holds
on $\mathcal B_k$. The union of the Gaussian bad events has
probability at most
\[
 C k^{260-5D+4D+192}=Ck^{452-D}.
\]
With $D=500$ this tends to zero. Choose the fixed grid-point
threshold $C_0k^{-D}$ large enough to account for the finitely many
coordinate norm comparisons. The interpolation error on
$\mathcal B_k$ is at most $Ck^{-D-2}$, so the grid bounds imply
$|u|+|\dd u|_{g_0}\ge k^{-D}W$ everywhere. Subtracting the bad
probabilities and \eqref{eq:coefficient-event} proves
\eqref{eq:noncancellation}. All order-dependent constants have only
increased the lower threshold on $k$.
\end{proof}

\subsection{Sign certificates with a power gain}

To obtain a fixed positive probability of opposite signs on each
wavelength ball, we need a different estimate. After normalizing
the variance of the value to one, we will show that at least one
directional derivative retains variance bounded away from zero even
when the value is fixed. The nonzero imaginary logarithmic gradients
in Lemma~\ref{lem:scaled-packets} provide this bound. Thus sign
production uses the packets' oscillation at the local wavelength,
while the preceding lemma uses their full first-jet span.

\begin{lemma}\label{lem:signs}
For every fixed choice of the packet orders in
\eqref{eq:gaussian-sum}, there exist $p_*,c_*>0$, independent of
$k$, such that for all sufficiently large allowed frequencies,
the probability that $u$ has a sign certificate in $\Omega$ of size
at least $c_*k^{1+\gamma}$ is at least $p_*$. The certificate can be
chosen from pairs of disks fixed deterministically before sampling
the Gaussian coefficients.
\end{lemma}

\begin{proof}
For $x\in V$, put
\begin{equation}\label{eq:sign-radius}
 R_x=(kQ_x)^{-1},\qquad Q_x=(1+|\nabla_G h(x)|_G^2)^{1/2}.
\end{equation}
All balls below are Euclidean coordinate balls. Their radii are at
most $k^{-1}$, so fixed enlargements lie in a common relatively
compact subset of $\Omega$ for large $k$. The envelope bounds give
$|DQ|\le Ck^{4\delta}$. If two candidate balls intersect, their
centers are at distance at most $2/k$, and their $Q$ values differ
by at most $Ck^{-1+4\delta}$. Since $Q\ge1$, the corresponding
radii are uniformly comparable.

At each fixed $k$ the candidate radii have a positive lower bound.
Choose a finite maximal disjoint family of open balls
$\{B(x,R_x):x\in\mathcal X_k\}$, with centers in $V$.
Uniformly bounded enlargements of these balls cover $V$, and $Q$ is
comparable to $Q_x$ in each enlargement. Thus
\begin{equation}\label{eq:weighted-packing}
 \int_V kQ(y)\,\dd y
 \le C\sum_{x\in\mathcal X_k}kQ_xR_x^4
 =C\sum_{x\in\mathcal X_k}R_x^3,
 \qquad
 \sum_{x\in\mathcal X_k}R_x^3\ge c k^{1+\gamma}.
\end{equation}
The last inequality uses \eqref{eq:envelope-bounds}.

Fix one ball. The beam vanishes on its radius-two enlargement for
large $k$, because $\overline V\subset\{t>1/2\}$ and $s_4<1/2$.
Define
\[
 S_x^2=\sum_{v\in\mathcal W_k}|v(x)|^2,
 \qquad U_x(Y)=u(x+R_xY)/S_x,\qquad |Y|\le2.
\]
In expressions differentiated with respect to $Y$, a packet $v$
means its composition $v(x+R_xY)$.
Lemma~\ref{lem:scaled-packets} applies to the entire packet
collection, including packets centered outside $V$ that meet this
ball. It gives $S_x\ge ce^{kh(x)}$; uniformly bounded sums of
normalized squared packet derivatives through every fixed order;
and, for every non-tail packet,
\begin{equation}\label{eq:local-imaginary-lower}
 \left|\operatorname{Im}\frac{\dd_Yv}{v}(0)\right|\ge c.
\end{equation}
Here a non-tail has its cutoff identically one on the scaled ball
and nonzero value at the center. For tails, the lemma supplies
absolute exponential bounds, even if $v(x)=0$. Polynomially many
tails consequently contribute $o(1)$ to the normalized squared
value and derivative sums. In particular,
\begin{equation}\label{eq:non-tail-mass}
 S_x^{-2}\sum_{v\ \mathrm{non\text{-}tail}}|v(x)|^2=1-o(1).
\end{equation}
All these statements are uniform over $x\in\mathcal X_k$.

The value $U_x(0)$ is a standard real Gaussian. We verify that some
directional derivative has a uniformly positive conditional
variance given this value. Let $z$ be the coefficient-space vector
whose coordinates for packet $v$ are
$(\operatorname{Re}v(x),\operatorname{Im}v(x))$, and let $T$ be the
matrix whose corresponding columns are the four-dimensional real
derivative vectors $\dd_Y\operatorname{Re}v(0)$ and
$\dd_Y\operatorname{Im}v(0)$. Thus $|z|^2=S_x^2$. If $g$ is the
vector of independent standard Gaussian coefficients, the value and
gradient are $z^{\mathsf t}g/S_x$ and $Tg/S_x$. The conditional
gradient covariance is
\begin{equation}\label{eq:conditional-covariance}
 S_x^{-2}T\left(I-\frac{zz^{\mathsf t}}{S_x^2}\right)T^{\mathsf t}.
\end{equation}
For each non-tail packet define the unit coefficient vector $c_v$
supported on its two coordinates, with entries
\[
 \frac{(-\operatorname{Im}v(x),\operatorname{Re}v(x))}{|v(x)|}.
\]
These vectors are mutually orthonormal and perpendicular to $z$.
The orthogonal projection in \eqref{eq:conditional-covariance}
therefore dominates $\sum_v c_vc_v^{\mathsf t}$ as a nonnegative
quadratic form. Moreover
\[
 Tc_v=|v(x)|\operatorname{Im}\frac{\dd_Yv}{v}(0).
\]
Equations~\eqref{eq:local-imaginary-lower} and
\eqref{eq:non-tail-mass} show that the trace of
\eqref{eq:conditional-covariance} is at least a fixed $c>0$.
There is consequently a deterministic unit direction $e_x$ in
$\R^4$ such that the derivative
$Z_x=\partial_{e_x}U_x(0)$, conditioned on $U_x(0)$, has variance
between fixed positive constants. Its unconditional variance is
bounded above by the deterministic squared-derivative estimates.
Those estimates also bound the regression coefficient
$\operatorname{Cov}(Z_x,U_x(0))$.

For clarity, the needed uniform control of the local $C^2$ norm
follows from finitely many derivatives. The deterministic estimates
give
\[
 \sup_{|Y|\le2}\E|D_Y^qU_x(Y)|^2\le C_q
 \qquad(0\le q\le5).
\]
Integration and the Sobolev inequality on concentric balls in
$\R^4$ yield
\begin{equation}\label{eq:local-c2-moment}
 \E\lVert U_x\rVert_{C^2(\overline{B(0,1)})}^2
 \le C\E\lVert U_x\rVert_{H^5(B(0,2))}^2\le C.
\end{equation}
We choose the $C^2$ norm to control the Euclidean gradient norm and
Hessian operator norm, changing only the fixed constant in this
inequality. The Sobolev embedding used here is valid since
$5>2+4/2$; see \cite[Sections~5.6--5.7]{Evans2010}.

Let $a'=1/4$ and fix $0<c'<a'/4$. For every sufficiently large
fixed number $H$, the Gaussian conditional density formula gives
\[
 \Pp\bigl(1\le Z_x\le2,\ |U_x(0)|<c'/H\bigr)\ge c_1/H.
\]
Indeed, on the indicated value interval the conditional mean is
bounded, and the conditional variance is bounded above and away
from zero. The Gaussian density on the fixed derivative interval
$[1,2]$ then has a uniform positive lower bound. By
\eqref{eq:local-c2-moment}, failure of the bound
$\lVert U_x\rVert_{C^2}\le H$ has probability at most $C/H^2$.
Choose $H$ once so large that the intersection of these three
conditions has probability at least
$\pi_*=c_1/(2H)>0$.

On that intersection Taylor's formula at $\pm a'e_x/H$ gives
opposite signs, with absolute values at least $a'/(2H)$:
the linear term has magnitude at least $a'/H$, the center value
has magnitude at most $c'/H$, and the quadratic error is at most
$a'^2/(2H)$. The gradient bound preserves these signs on the
parallel transverse disks of radius
\[
 \sigma_* = a'/(4H^2)
\]
centered at those two points. The whole closed cylinder between
the disks is compactly contained in the unit ball. The direction
$e_x$, centers and radii are deterministic functions of the packet
data. After scaling by $R_x$, they provide a potential certificate
pair inside $B(x,R_x)$.

Let $F_k$ be the fraction of successful balls weighted by $R_x^3$.
No independence among these successes is claimed or needed.
Since $0\le F_k\le1$ and $\E F_k\ge\pi_*$,
\begin{equation}\label{eq:weighted-success}
 \Pp(F_k\ge\pi_*/2)\ge\frac{\pi_*}{2-\pi_*}>0.
\end{equation}
The Euclidean three-area of one successful disk is
$\omega_3\sigma_*^3R_x^3$, where $\omega_3$ is the unit-ball
volume in $\R^3$. On the event in \eqref{eq:weighted-success},
\eqref{eq:weighted-packing} therefore gives a certificate of size
at least $c_*k^{1+\gamma}$. Its cylinders lie in disjoint open
balls, as required by Definition~\ref{def:certificate}. This proves
the lemma with $p_* =\pi_* /(2-\pi_*)$.
\end{proof}

Lemmas~\ref{lem:noncancellation} and~\ref{lem:signs} can be realized
simultaneously. For each fixed choice of the orders, first fix the
positive number $p_*$ in Lemma~\ref{lem:signs}, and then take $k$
so large that the failure probability in
\eqref{eq:noncancellation} is less than $p_*/2$. The simultaneous
event has probability at least $p_*/2$. Thus, for every sufficiently
large allowed $k$, there is a real smooth $u$ satisfying
\eqref{eq:sum-weighted-bounds}, the pointwise lower bound in
\eqref{eq:noncancellation}, and a certificate of size
$c_*k^{1+\gamma}$. This probability subtraction is valid even if
$p_*$ is extremely small. The probability bounds above were obtained
without conditioning on $\mathcal B_k$; only deterministic derivative
estimates and interpolation used that event.

\subsection{Exact coefficient correction}

\begin{lemma}\label{lem:correction}
Let $(A,\rho)$ be a fixed smooth positive pair, let $u$ be real and
smooth, and let $W$ be a smooth positive weight. Suppose that
\[
 |u|+|\dd u|_{g_0}\ge k^{-500}W.
\]
Fix an integer $R\ge0$. Suppose that
\[
 |D^qu|\le C_qk^{2q+41}W\quad(0\le q\le R+2),\qquad
 |D^q\mathfrak r|\le C_qk^{-M}W\quad(0\le q\le R+1),
\]
where $\mathfrak r=\diver_\mu(A\,\dd u)+k^2\rho u$.
There are smooth real corrections $(\delta A,\delta\rho)$ canceling
the eigenfunction residual exactly, and
\begin{equation}\label{eq:correction-size}
 \lVert(\delta A,\delta\rho)\rVert_{C^R}
 \le C_R k^{-M+2129+1088R}.
\end{equation}
The corrections vanish on every open set where $\mathfrak r=0$.
In particular, choosing $M\ge2131+1088R$ makes them tend to zero
in $C^R$ as $k\to\infty$. For sufficiently small corrections the
output pair is positive and $u$ is its exact eigenfunction at $k^2$.
\end{lemma}

\begin{proof}
Write $I=g_0^{-1}$ and define globally
\begin{equation}\label{eq:correction-formula}
 \begin{split}
 \mathcal D&=u^2+I(\dd u,\dd u),\qquad F=-\mathfrak r/\mathcal D,\\
 \delta A&=FuI,\qquad
 \delta\rho=k^{-2}\bigl(Fu-\diver_\mu(FI\,\dd u)\bigr).
 \end{split}
\end{equation}
The lower bound on the first jet gives
$\mathcal D\ge\tfrac12 k^{-1000}W^2>0$, so these are smooth.
The product rule gives the exact identity
\begin{align*}
 \diver_\mu(\delta A\,\dd u)+k^2\delta\rho\,u
 &=u\diver_\mu(FI\,\dd u)+F I(\dd u,\dd u)\\
 &\quad+Fu^2-u\diver_\mu(FI\,\dd u)\\
 &=F\mathcal D=-\mathfrak r.
\end{align*}
Thus $\diver_\mu((A+\delta A)\dd u)
+k^2(\rho+\delta\rho)u=0$ identically.

We include finite derivative bounds to specify the accuracy needed
in this step. For $0\le q\le R+1$, the hypotheses imply
\[
 |D^q\mathcal D|\le C_q k^{2q+86}W^2.
\]
Every term in a derivative of order $q$ of $\mathcal D^{-1}$ is a
constant times
$\mathcal D^{-m-1}\prod_{i=1}^m D^{q_i}\mathcal D$, where
$q_i\ge1$, $\sum_iq_i=q$, and $m\le q$; the order-zero case is
included separately. Hence
\[
 |D^q(\mathcal D^{-1})|
 \le C_q k^{1000+2q+1086q}W^{-2}
 =C_q k^{1000+1088q}W^{-2}.
\]
By the product rule,
\begin{equation}\label{eq:F-derivatives}
 |D^qF|\le C_qk^{-M+1000+1088q}W^{-1}
 \qquad(0\le q\le R+1).
\end{equation}
No derivatives of $W^{-1}$ are required for these estimates: the
displayed weights bound each differentiated factor before the
factors are multiplied.

Multiplying \eqref{eq:F-derivatives} by derivatives of $u$ gives
$|D^q(Fu)|\le C_qk^{-M+1041+1088q}$. Replacing $u$ by $I\dd u$
gives the bound $C_qk^{-M+1043+1088q}$, with the derivatives of
the fixed tensor $I$ absorbed in the constant. The divergence in
\eqref{eq:correction-formula} costs one further derivative, and
the density correction has the factor $k^{-2}$. These bounds give
\eqref{eq:correction-size}; its exponent for the density is
$1043+1088(R+1)-2=2129+1088R$.
The cancellation of the two weights is what permits arbitrarily
small coefficient changes even where the original function is
exponentially small.

If $\mathfrak r=0$ on an open set, then $F=0$ there and both
corrections vanish. Finally, the cone of positive pairs is open
in $C^0$ on the compact base, so sufficiently small corrections
preserve positivity.
\end{proof}

\subsection{Simplicity on a protected round patch}

\begin{lemma}\label{lem:simplicity}
Suppose a smooth positive pair has a real eigenfunction $u$ at
$k^2=N(N+3)$, with $N>1$. Suppose there is an open region on which
the pair equals the round pair and $u=u_*$, containing a point with
$0<t<1$ at which $\dd u$
and $\dd t$ are independent. Then one can make $k^2$ simple by
an arbitrarily small smooth tensor perturbation supported in that
region, while preserving $u$ and $\rho$ exactly. The perturbation
can be arbitrarily small in any prescribed finite $C^R$ norm.
\end{lemma}

\begin{proof}
Take a small product coordinate patch $O$ in the stated region on
which $(u,t)$ are the first two coordinates, and $\dd u\ne0$.
Choose a smooth nonnegative bump $\chi$ supported in $O$ and
strictly positive on a smaller product patch $O_0$. Define
\[
 E=\chi\left(I-
 \frac{(I\dd u)\otimes(I\dd u)}{I(\dd u,\dd u)}\right),
 \qquad I=g_0^{-1},
\]
extending by zero off $O$. This is a smooth nonnegative symmetric
contravariant tensor and $E\dd u=0$.

Let $\mathcal E$ be the real eigenspace at $k^2$ for the given
pair. The nonnegative form
\[
 f\longmapsto\int_B E(\dd f,\dd f)\,\dd\mu,
 \qquad f\in\mathcal E,
\]
has kernel exactly $\operatorname{span}\{u\}$. Indeed, a kernel
element satisfies $\dd f\in\operatorname{span}\{\dd u\}$ on
$O_0$. After making the product patch smaller, this means
$f=b(u)$ for a smooth one-variable function $b$. The round
eigenfunction equations imply
\begin{equation}\label{eq:simplicity-composition}
 b''(u)|\dd u|_{g_0}^2=k^2\bigl(ub'(u)-b(u)\bigr).
\end{equation}
The Euclidean gradient of the homogeneous harmonic polynomial
defining $u_*$ has squared norm $N^2(1-t)^{N-1}$, and its radial
derivative is $Nu$. Tangential projection onto $S^4$ consequently
gives
\[
 |\dd u|_{g_0}^2=N^2\bigl((1-t)^{N-1}-u^2\bigr).
\]
Differentiate \eqref{eq:simplicity-composition} with respect to $t$
while fixing the coordinate $u$. Since $N>1$ and $t<1$, it follows
that $b''(u)=0$. Substitution back into
\eqref{eq:simplicity-composition} eliminates the constant term of
$b$, so $f=cu$ on $O_0$.

The difference $f-cu$ satisfies a second-order elliptic equation
with smooth real positive definite principal coefficients and
smooth lower-order coefficients. Aronszajn's unique continuation
theorem applies: its local hypotheses allow principal coefficients
of class $C^2$ with Lipschitz second derivatives and bounded
lower-order coefficients, and smooth solutions satisfy its
solution hypotheses \cite[pp.~235--236]{Aronszajn1957}. Vanishing
on an open patch gives infinite-order vanishing there. On the
connected base the difference is therefore identically zero,
proving the assertion about the kernel.

For $\tau>0$ replace $A$ by $A+\tau E$. The eigenfunction $u$ and
its eigenvalue remain exact because $E\dd u=0$ identically. We
claim that $k^2$ is simple for every sufficiently small positive
$\tau$. Otherwise there would be $\tau_m\downarrow0$ and real
eigenfunctions $f_m$ at $k^2$, normalized in $L^2(\rho\,\dd\mu)$
and orthogonal to $u$. Their energy identity is
\[
 \int_B(A+\tau_m E)(\dd f_m,\dd f_m)\,\dd\mu=k^2,
\]
so uniform ellipticity bounds them in $H^1$. A subsequence
converges weakly in $H^1$ and strongly in $L^2$ to a normalized
function $f_0$ orthogonal to $u$. Passing to the weak equation
shows $f_0\in\mathcal E$; elliptic regularity makes it smooth.
For any fixed $\phi\in\mathcal E$, testing the two eigenfunction
equations against one another cancels the original form and gives
\[
 \int_B E(\dd f_m,\dd\phi)\,\dd\mu=0.
\]
Weak convergence passes this fixed bilinear identity to the limit.
Now choose $\phi=f_0$ to see that $f_0$ is in the kernel just
identified, contradicting its normalization and orthogonality.
This proves the claim. The tensor $E$ is fixed before choosing
$\tau$; hence $\tau E$ can meet every prescribed finite norm
tolerance, while remaining supported in $O$.
\end{proof}

\subsection{Proof of the quantitative step}

\begin{proof}[Proof of Proposition~\ref{prop:step}]
Fix the input pair, the collar, $r$ and $\eta$. Choose
$R\ge\max\{r,1\}$ and a smaller positive tolerance $\eta'$ so
that total $C^R$ movement less than $\eta'$ implies the requested
$C^r$ movement and keeps the pair positive and in $\mathcal U$.
Such a choice exists because the fixed input belongs to the open
set $\mathcal U$.

Choose $J=R+2$ and $M\ge2131+1088R$, and then choose the finite
packet orders giving \eqref{eq:sum-weighted-bounds}. All these
choices precede frequency selection. Lemmas~\ref{lem:noncancellation}
and~\ref{lem:signs} give, for every sufficiently large allowed $k$,
a sample with the required first-jet bound, residual bounds and a
certificate of size $c_*k^{1+\gamma}$. Lemma~\ref{lem:correction}
changes the pair by at most $C_Rk^{-2}$ in $C^R$ while preserving
that same function. Choose $k$ large enough for this change to be
less than $\eta'/2$.

We describe explicitly why a narrow collar only changes this
frequency threshold. Packet supports avoid the region below
$s_1-Ck^{-1/3-100\delta}$. Requiring
\begin{equation}\label{eq:collar-thresholds}
 Ck^{-1/3-100\delta}<w/100,
 \qquad
 cw k>(M'+P)\log k+\log C_q(w)
\end{equation}
protects a fixed subcollar and makes any required cutoff residual
$C_q(w)k^Pe^{-cw k}$ smaller than $k^{-M'}$. Only finitely many
orders $q$ and finite targets $M'$ are needed at this step.
For every already fixed $w>0$, all inequalities in
\eqref{eq:collar-thresholds} hold for every sufficiently large $k$.
In particular the exact corrections vanish below $s_1-w/100$,
where the coefficients remain round and $u=u_*$.

In the subcollar
\[
 l_-+w/20<t<l_-+w/10
\]
choose a point with $\sin(N\vartheta)\ne0$. Since $0<t<1$,
$\dd t\ne0$, and the nonzero angular derivative makes $\dd u$
and $\dd t$ independent. Lemma~\ref{lem:simplicity} applies there.
After $k$ and the sample are fixed, form its tensor $E$ and choose
$\tau>0$ below its simplicity threshold and so small that
$\lVert\tau E\rVert_{C^R}<\eta'/2$. This tensor can have large
finite derivatives depending on $k$; the subsequent scalar choice
of $\tau$ absorbs them. It preserves $u$, its exact eigenvalue and
all current signs exactly. Its support lies in $t>l_-$, so the
final pair is round on $t<l_-$.

The final pair meets every conclusion of
Proposition~\ref{prop:step}. In particular, its certificate constant
$c_*$ was fixed after choosing the finite orders and before taking
$k$ large, and it was unaffected by either correction. The
exponent $\gamma$, the chart and the background neighborhood were
fixed upstream. Neither the collar estimates nor the arbitrary
finite tolerances change them. The construction works for every
sufficiently large integer $N$, as asserted.
\end{proof}

\subsection{Spectral persistence under finite norm tolerances}

\begin{lemma}\label{lem:persistence}
Let a smooth positive pair $\mathcal P=(A,\rho)$ have a simple
positive base eigenvalue $\lambda$ and a real eigenfunction with a
sign certificate $\mathcal C$. If $I\subset(0,\infty)$ is an open
interval containing $\lambda$, there are an integer $s\ge1$ and
$\varepsilon>0$ such that every smooth positive pair
$\mathcal P'$ satisfying
\[
 \lVert\mathcal P'-\mathcal P\rVert_{C^s}\le\varepsilon
\]
has a simple eigenvalue in $I$ and a real eigenfunction with the
same certificate disks. Thus a closed ball defined by one finite
norm can be contained in a persistence neighborhood.
\end{lemma}

\begin{proof}
We recall the spectral and regularity facts used here. The form
\[
 \int_B A(\dd f,\dd v)\,\dd\mu+
 \int_B\rho fv\,\dd\mu
\]
is coercive on $H^1(B)$ and defines the inverse of $-L+1$ in
$L^2(\rho\,\dd\mu)$. The compact embedding
$H^1\hookrightarrow L^2$ makes this inverse compact, and it is
selfadjoint. The spectral theorem and elliptic regularity give a
discrete nonnegative spectrum with finite multiplicities and smooth
real eigenfunctions. Its ordered eigenvalues obey the min--max
principle with Rayleigh quotient
\begin{equation}\label{eq:weighted-rayleigh}
 \frac{\int_B A(\dd f,\dd f)\,\dd\mu}
      {\int_B\rho f^2\,\dd\mu}.
\end{equation}

Uniform multiplicative comparisons of nearby positive tensors and
densities compare \eqref{eq:weighted-rayleigh} for every $f$ and
therefore imply convergence of each ordered eigenvalue under
$C^0$ convergence of pairs. Choose a smaller interval about
$\lambda$, with closure in $I$, separating it from its neighboring
ordered eigenvalues. Their convergence implies that for sufficiently
close pairs the corresponding ordered eigenvalue is still simple
and lies in $I$.

We also need uniform convergence of a corresponding normalized
eigenfunction, up to sign. Here only finitely many coefficient
derivatives are needed. For a fixed smooth elliptic $L$ on the
closed base, the local interior estimates, a finite partition of
unity and interpolation give
\begin{equation}\label{eq:elliptic-persistence-estimate}
 \lVert f\rVert_{H^{a+2}}
 \le C_a\bigl(\lVert Lf\rVert_{H^a}+\lVert f\rVert_{L^2}\bigr)
 \qquad(a=0,1,2,\ldots).
\end{equation}
These are the usual interior elliptic estimates; the Sobolev and
compactness statements used here are described in
\cite[Sections~5.6--5.7 and~6.3]{Evans2010}. A sufficiently small
finite coefficient norm makes $L'-L$ small as an operator
$H^{a+2}\to H^a$ for each of finitely many specified $a$. Its
contribution in \eqref{eq:elliptic-persistence-estimate} can then
be absorbed on the left. Applying the resulting estimates at
$a=0,2$, together with the eigenfunction equation, gives a uniform
$H^4$ bound for eigenfunctions whose eigenvalues stay in the
chosen bounded interval and whose $L^2$ norms are bounded.
For example, requiring sufficiently small $C^4$ changes of the
pair is more than enough for these finitely many operator bounds:
the principal coefficients of $L$ involve $A,\rho$ without
derivatives, and its first-order coefficients involve at most
their first derivatives.

Suppose that pairs converge in such a finite norm to $\mathcal P$.
Normalize their corresponding real eigenfunctions $f_m$ by
$\int\rho_m f_m^2\,\dd\mu=1$. Positivity bounds their ordinary
$L^2$ norms, so the preceding argument bounds them in $H^4$.
Rellich compactness and Sobolev embedding in dimension four give
a subsequence converging in $C^0$ (for example, first converge in
$H^3$, which embeds in $C^0$). Passing to the weak equation shows
that the limit is a normalized eigenfunction at $\lambda$ for
$\mathcal P$. By simplicity it equals the prescribed normalized
eigenfunction or its negative.

The finite collection of strict signs on compact disks has a
positive minimum absolute margin after this normalization. Hence
all sufficiently close such eigenfunctions have the same opposite
signs on the same disks, allowing an overall change of sign. If
no sufficiently small finite-norm neighborhood had this property,
a sequence converging in that norm would contradict the preceding
compactness argument. Finally, decrease its radius to put a closed
finite-norm ball inside the neighborhood. This proves the lemma,
including the assertion for varying densities.
\end{proof}

\subsection{One smooth limiting pair and the ordinary metric}

\begin{proof}[Proof of Theorem~\ref{thm:main}]
Use monotone fixed coordinate norms for pairs. Choose $b>0$ so
that the closed $C^1$ ball of radius $b$ about
$\mathcal P_0=(g_0^{-1},1)$ consists of positive pairs in
$\mathcal U$. Fix
\[
 l^{(j)}=\frac14+2^{-j-2}\quad(j\ge0),\qquad
 \epsilon_0=\gamma/4>0.
\]
We construct pairs $\mathcal P_j=(A_j,\rho_j)$, frequencies $k_j$,
and certificates $\mathcal C_j$ inductively. Every output is round
on $t<l^{(j)}$.

After producing $\mathcal P_i$, apply Lemma~\ref{lem:persistence}
to its simple eigenvalue $k_i^2$, its certificate, and the interval
$(k_i^2/2,2k_i^2)$. Denote a resulting finite norm order and
closed-ball radius by $s_i$ and $\varepsilon_i>0$.
At stage $j$, these numbers are already fixed for $i<j$. Put
\begin{equation}\label{eq:diagonal-budgets}
 \begin{split}
 r_j&=\max\{j+1,1,s_1,\ldots,s_{j-1}\},\\
 \eta_j&=\min\left\{2^{-j},\ b2^{-j-1},\,
 \min_{1\le i<j}2^{-(j-i+1)}\varepsilon_i\right\}.
 \end{split}
\end{equation}
For $j=1$, the inner minimum is omitted. These are finite
requirements at each stage. Apply Proposition~\ref{prop:step}
to $\mathcal P_{j-1}$, with
$l_+=l^{(j-1)}$, $l_-=l^{(j)}$, order $r_j$ and tolerance
$\eta_j$. Use the certificates from its construction in
Lemma~\ref{lem:signs}, with balls centered in $V$ and radii at most
$k_j^{-1}$. The proposition gives a positive constant $c_j$ and all
sufficiently large allowed frequencies. Choose one with
$k_j>2k_{j-1}$ when $j>1$, and also so large that
\[
 c_jk_j^{\gamma/2}>j.
\]
The output certificate then satisfies
\begin{equation}\label{eq:stage-certificate-size}
 |\mathcal C_j|>j k_j^{1+\gamma/2}.
\end{equation}
Only after this output has been constructed do we select its own
$s_j,\varepsilon_j$ by Lemma~\ref{lem:persistence}. They constrain
subsequent stages, so no current frequency choice depends
circularly on its own future persistence radius.

The increments are summable in every fixed $C^r$ norm: for all
large $j$, $r_j\ge r$ and their $C^{r_j}$ norms are less than
$2^{-j}$. Therefore the pairs converge to one smooth pair
$\mathcal P_\infty=(A_\infty,\rho_\infty)$. Their total $C^1$
movement from the round pair is less than
$b\sum_{j\ge1}2^{-j-1}=b/2$, so the limit remains positive and
belongs to $\mathcal U$. For each fixed $i$, the same budgets give
\[
 \lVert\mathcal P_\infty-\mathcal P_i\rVert_{C^{s_i}}
 \le\sum_{j>i}2^{-(j-i+1)}\varepsilon_i
 =\varepsilon_i/2.
\]
The limit lies in every selected persistence neighborhood. Hence
there are real smooth eigenfunctions $v_i$ for the limiting base
operator, with the same certificate disks $\mathcal C_i$, and
\begin{equation}\label{eq:limiting-eigenvalue-intervals}
 -L_\infty v_i=\lambda_i v_i,\qquad
 k_i^2/2<\lambda_i<2k_i^2.
\end{equation}
The frequency separation makes these eigenvalue intervals disjoint
and tending to infinity. The functions $v_i$ are nonzero by their
normalization, or just by their strict signs.

Set $G_\infty^{-1}=A_\infty/\rho_\infty$ and let
\[
 q=\frac{\rho_\infty\,\dd\mu}{\dd\operatorname{vol}_{G_\infty}}.
\]
This is a smooth positive function, the ratio of two smooth
positive base densities. We use the standard warped-product
realization of a weighted Laplacian
\cite[Section~4, equation~(2)]{CharalambousLu2015}, with a circle
as the fiber. On the fixed manifold $M=B\times S^1$
define
\begin{equation}\label{eq:circle-metric}
 g=G_\infty+q^2\dd\varphi^2,
\end{equation}
where $\dd\varphi^2$ is the standard unit-circle metric. The
manifold is smooth, closed and connected, of dimension five, and
\[
 \dd\operatorname{vol}_g
 =q\,\dd\operatorname{vol}_{G_\infty}\,\dd\varphi
 =\rho_\infty\,\dd\mu\,\dd\varphi.
\]
For a function $f$ pulled back from $B$, the coordinate divergence
formula gives
\begin{equation}\label{eq:circle-operator}
 \Delta_g f
 =\frac{1}{\rho_\infty\mu}
   \partial_a\bigl(\rho_\infty\mu
                    G_\infty^{ab}\partial_b f\bigr)
 =\rho_\infty^{-1}\diver_\mu(A_\infty\dd f)
 =L_\infty f.
\end{equation}
Here $\mu$ in the coordinate expression denotes its positive
coordinate density. Thus the pullbacks $u_i$ of $v_i$ are exact
real Laplace--Beltrami eigenfunctions with eigenvalues
$\lambda_i$.

It remains to convert their certificates into the claimed measure
bound. Fix a closed interval $I_0$ of positive length in a circle
coordinate chart. All certificate cylinders are in one fixed
compact subset of $\Omega$: their centers lie in $V\Subset\Omega$
and their radii tend uniformly to zero, with the initial
frequencies chosen large enough for the same inclusion. On the
product of this compact set with $I_0$, coordinate and Riemannian
distances have fixed Lipschitz comparison constants. One may see
this directly by extending the coordinate functions smoothly with
cutoffs to the compact manifold; their gradients are then
uniformly bounded in the single fixed metric $g$.

For one certificate cylinder $C_i^{(a)}$, let $D_i^{(a)}$ be one
of its transverse disks. The opposite endpoint signs imply that
every longitudinal segment through the disk contains a zero of
$v_i$. Therefore the subset
\[
 \Z(u_i)\cap(C_i^{(a)}\times I_0)
\]
projects onto $D_i^{(a)}\times I_0$ by coordinate orthogonal
projection in the base and the circle coordinate. The Lipschitz
constant of every such projection is bounded by one fixed $K$,
independent of the disk's radius and orientation. The defining
covering inequality for Hausdorff measure under a Lipschitz map
gives
\[
 \Hh_g^4\bigl(\Z(u_i)\cap(C_i^{(a)}\times I_0)\bigr)
 \ge K^{-4}\Hh_{\mathrm{Euc}}^4(D_i^{(a)}\times I_0)
 =K^{-4}|I_0|\Hh_{\mathrm{Euc}}^3(D_i^{(a)}).
\]
The closed cylinders lie in disjoint open balls. Their indicated
nodal subsets are consequently pairwise disjoint Borel sets, so
their Hausdorff measures add. Using
\eqref{eq:stage-certificate-size}, we obtain one constant $c>0$
belonging to the final metric such that
\begin{equation}\label{eq:final-nodal-lower}
 \Hh_g^4(\Z(u_i))\ge c|\mathcal C_i|
 > c i k_i^{1+\gamma/2}.
\end{equation}
This argument uses continuity and signs only, without any
regularity assumption on the nodal set.

Finally, \eqref{eq:limiting-eigenvalue-intervals} and
$\epsilon_0=\gamma/4$ imply
\[
 \lambda_i^{1/2+\epsilon_0}
 \le2^{1/2+\epsilon_0}k_i^{1+\gamma/2}.
\]
Combining this with \eqref{eq:final-nodal-lower} yields
\[
 \frac{\Hh_g^4(\Z(u_i))}{\lambda_i^{1/2+\epsilon_0}}
 \ge c\,2^{-1/2-\epsilon_0}i\longrightarrow\infty.
\]
Both $g$ and $\epsilon_0$ are fixed throughout this sequence,
which proves Theorem~\ref{thm:main}.
\end{proof}

\begin{corollary}
For every integer $n>5$ there is a smooth closed connected
$n$-dimensional Riemannian manifold and a fixed $\epsilon_0>0$
with the same divergence of nodal measure divided by
$\lambda^{1/2+\epsilon_0}$.
\end{corollary}

\begin{proof}
Take the product of \eqref{eq:circle-metric} with any fixed smooth
closed connected manifold of dimension $n-5$, and pull back the
same eigenfunctions. The product Laplacian preserves their
eigenvalues. Multiply each of the projection certificates above
by a fixed compact coordinate box of positive $(n-5)$-volume in
the extra factor. The same Lipschitz projection argument, now in
dimension $n-1$, gives a fixed positive multiple of the preceding
certificate lower bound. The denominator and its fixed exponent
are unchanged, so the ratio still tends to infinity.
\end{proof}

\bibliographystyle{plain}
\bibliography{references}
\end{document}
